% !TeX root = ../../Book3.tex
% 纲要标题：## 第9章　截面曲率、空间形式与三角比较
% 纲要位置：几何学计划纲要.md，第 2426 行。
% 正文状态：初稿；尚未正式审核。

% 纲要内容开始
% 
% > 本章目标与证明范围：把上界和下界的不同三角几何分开。Cartan–Hadamard 与空间形式分类给完整证明；Toponogov 的选定全局版本有独立的局部到整体论证。
% 
% 纲要内容结束

\chapter{截面曲率、空间形式与三角比较}
\label{g3:ch:09}
\BookChapterNumber{g3:ch:09}

\begin{chapterabstract}
	截面曲率的上界和下界怎样约束整个测地三角形？本章先在非正曲率下把无共轭点与完备性结合，证明指数映射是万有覆盖，进而在单连通情形得到全局唯一测地线和距离凸性。常曲率 Jacobi 方程随后确定局部模型；对正曲率极坐标的公共终点作光滑延拓，完成单连通模型及离散群商的整体分类。Toponogov 比较通过一维上支撑拼接，给出严格模型周长范围内的全局顶点到对边及角度比较。最后建立闭凸集投影与 CAT(0) 性质，辨明同初始角铰链和同边长三角形所使用的数据与比较方向。
\end{chapterabstract}

\begin{learninggoals}
	\begin{itemize}
		\item 证明非正曲率下指数映射的微分扩张，并说明完备性和单连通性如何分别进入整体结论。
		\item 由 Jacobi 场恢复常曲率局部度量，处理正曲率模型的对径点并完成空间形式分类。
		\item 区分同边长三角形和同边同角铰链，追踪 Toponogov 的有效周长及割点处理。
		\item 证明 Hadamard 流形中闭凸集的最近点投影存在、唯一且不增加距离。
		\item 用有限点的距离比较表达上、下曲率条件，并用平坦商检验局部与整体的差别。
	\end{itemize}
\end{learninggoals}

在 Euclid 平面上，起点与初速度确定一条直线，任意两点之间又恰有一条最短线。这两个事实在一般黎曼流形上分属不同层面：测地线初值问题仍有唯一解，但圆球面的对径点有许多最短连接，平坦环面的不同绕行也可能到达同一终点。曲率控制指数映射的局部微分，覆盖空间则记录全局重复。要理解三角形，必须同时看见这两部分。

\ProseParagraphBreak

第七章比较了径向 Jacobi 场，第八章对方向求和、对半径积分，得到了 Ricci 曲率所控制的直径与体积。本章保留逐截面的信息。非正曲率使拉回到切空间上的度量不小于 Euclid 度量，这个长度下界保证拉回度量完备，从而使局部等距的指数映射成为覆盖。若曲率恒定，Jacobi 方程还能精确求解；剩下的问题是怎样将局部模型拼成整体，以及哪些离散等距群可以产生不同的商。

\ProseParagraphBreak

三角比较的另一个难点是比较哪些数据。固定两条边及夹角时，较大的曲率使对边受到压缩；固定三条边时，相应结论变为角度及边内点距离的另一方向。正曲率模型还存在首个共轭尺度，故周长条件参与模型三角形的存在与唯一性。我们在\secref{g3:sec:0903}中证明所选的严格周长版本，并在使用上支撑的一维引理中完成从局部估计到整条对边的拼接。

\ProseParagraphBreak

所需基础是第二章的指数映射与完备性、第五章的完备局部等距覆盖引理、第六章的 Jacobi 场及变分公式、第七章的 Hessian 比较和移动基点支撑。常曲率及非正曲率的阅读对照可查 Petersen 第三版第5.6.2节、第6.2.1--6.2.2节，三角比较可查第12.2节\cite[第5.6.2节，第6.2.1--6.2.2节，第12.2节]{Petersen2016}。最后一节的投影和 CAT(0) 证明不依赖选读节；下曲率三角比较则明确回用该节的结论。

% !TeX root = ../../Book3.tex
% 纲要标题：### 9.1 非正曲率与 Cartan–Hadamard 定理
% 纲要位置：几何学计划纲要.md，第 2430 行。
% 正文状态：初稿；尚未正式审核。

% 纲要内容开始
% 
% * 无共轭点与指数映射
% * 覆盖论证及单连通条件
% * 全局唯一测地线
% * 距离函数凸性和拓扑后果
% 
% 纲要内容结束

\section{非正曲率与 Cartan–Hadamard 定理}
\label{g3:sec:0901}

第七章的比较定理说明，非正截面曲率使同一点发出的径向 Jacobi 场至少按 Euclid 模型的速度分离，因而指数映射没有奇异点。然而，平坦环面上不同的初速度仍可能到达同一点：微分可逆性控制局部，单连通性才排除绕行造成的重复。本节把这两方面结合起来，证明 Cartan--Hadamard 定理，并把局部的 Hessian 比较延伸为全局凸性。\footnote{Jacques Salomon Hadamard（1865--1963），法国数学家，研究解析数论、偏微分方程及负曲率曲面上的测地线等问题；他于1896年证明了素数定理。}

\subsection{无共轭点与指数映射}
\label{g3:subsec:0901:01}

先把曲率条件能够直接给出的结论分离出来。这里尚不假定单连通；即使指数映射最终不是单射，下面的微分估计仍然成立。

\begin{proposition}\label{g3:prop:nonpositive-exponential-differential}
	设 \((M^n,g)\) 是连通、完备、光滑无边界黎曼流形，\(n\ge1\)，且所有截面曲率满足 \(K\le0\)；当 \(n=1\) 时，截面条件为空。则任意测地线上都没有共轭点。对每个 \(p\in M\)，指数映射定义在整个 \(\mathrm{T}_pM\) 上，并且
	\begin{equation}\label{g3:eq:nonpositive-exponential-expansion}
		|(\mathrm d\exp_p)_v(\xi)|_g\ge|\xi|_{g_p},
		\qquad v,\xi\in\mathrm{T}_pM.
	\end{equation}
	这里用向量空间的平移识别 \(\mathrm{T}_v(\mathrm{T}_pM)\) 与 \(\mathrm{T}_pM\)。特别地，\(\exp_p\) 在每点都是局部微分同胚。
\end{proposition}
\begin{proof}
	沿任意非恒定测地线 \(\gamma\)，记 \(T=\dot\gamma\)。若 Jacobi 场 \(J\) 在两个参数值 \(a<b\) 消失，则 Jacobi 方程给出
	\[
		\frac{\mathrm d^2}{\mathrm dt^2}|J|^2
		=2|J'|^2-2\langle R(J,T)T,J\rangle\ge0.
	\]
	最后一步来自 \(K\le0\)：只有 \(J\) 垂直于 \(T\) 的分量参与曲率项，且本书约定下正曲率对应非负的 \(\langle R(J,T)T,J\rangle\)。一维时曲率张量为零，同一不等式成立。非负函数 \(|J|^2\) 凸且在 \(a,b\) 处为零，故在整个 \([a,b]\) 上为零；Jacobi 初值唯一性随后使 \(J\) 恒为零。这排除了共轭点。

	完备性与 \cref{g3:thm:hopf-rinow}保证 \(\exp_p\) 在整个切空间上有定义。对 \(v\ne0\)，写 \(v=ru\)、\(|u|=1\)，并把 \(\xi\) 分解为
	\[
		\xi=\alpha u+w,\qquad w\perp u.
	\]
	沿单位速径向线 \(\gamma(t)=\exp_p(tu)\)，令 \(J_w(0)=0\)、\(J_w'(0)=w\)。指数映射的 Jacobi 表达和径向、横向的正交性给出
	\[
		(\mathrm d\exp_p)_{ru}(u)=T(r),\qquad
		(\mathrm d\exp_p)_{ru}(w)=\frac{J_w(r)}r,\qquad
		|(\mathrm d\exp_p)_{ru}\xi|^2
		=\alpha^2+\frac{|J_w(r)|^2}{r^2}.
	\]
	当 \(n\ge2\) 时，\cref{g3:thm:rauch-model-comparison}取 \(k=0\)，模型函数为 \(s_0(t)=t\)，且模型没有正零点，所以 \(|J_w(r)|\ge r|w|\) 对所有 \(r>0\) 成立。这给出式\eqref{g3:eq:nonpositive-exponential-expansion}。一维时 \(w=0\)，结论直接由径向等式得到。在 \(v=0\) 处，\((\mathrm d\exp_p)_0=\operatorname{id}\)。

	该不等式使指数映射的微分为单射，源与靶同维，因而微分可逆；逆函数定理给出局部微分同胚性。
\end{proof}

从无共轭点也可利用 \cref{g3:prop:conjugate-exp-singularity}直接得到微分可逆性。上面的长度估计还多提供了一个整体论证所需的信息：在切空间上拉回度量以后，向量的长度不会小于原来的 Euclid 长度。

\subsection{覆盖论证及单连通条件}
\label{g3:subsec:0901:02}

固定 \(p\in M\)，在向量空间 \(V=\mathrm{T}_pM\) 上定义
\begin{equation}\label{g3:eq:exponential-pullback-metric}
	h_v(\xi,\eta)
	=g_{\exp_p(v)}\bigl((\mathrm d\exp_p)_v\xi,
		(\mathrm d\exp_p)_v\eta\bigr).
\end{equation}
微分处处可逆，所以 \(h=\exp_p^*g\) 是光滑正定的黎曼度量。按定义，\(\exp_p:(V,h)\to(M,g)\) 是局部等距映射；而式\eqref{g3:eq:nonpositive-exponential-expansion}表明
\[
	h_v(\xi,\xi)\ge g_p(\xi,\xi).
\]
记 \(g_0\) 为 \(V\) 上由内积 \(g_p\) 给出的常系数 Euclid 度量。下面证明 \(h\ge g_0\) 在这里确实保证完备。

\begin{lemma}\label{g3:lem:euclidean-dominated-metric-complete}
	设 \(V\) 是有限维实内积空间，\(g_0\) 为其 Euclid 度量，\(h\) 是定义在整个 \(V\) 上的光滑黎曼度量。若 \(h\ge g_0\)，则 \((V,h)\) 度量完备。
\end{lemma}
\begin{proof}
	对任意分片光滑曲线 \(c\)，有 \(L_h(c)\ge L_{g_0}(c)\)，取下确界得到
	\[
		d_h(v,w)\ge |v-w|_{g_0}.
	\]
	因此每个 \(d_h\)-Cauchy 序列 \((v_j)\) 都是 Euclid 距离的 Cauchy 序列，收敛于某个 \(v_\infty\in V\)。

	还须验证它在 \(d_h\) 中也收敛。取以 \(v_\infty\) 为中心的 Euclid 闭球 \(K\)，半径为某个 \(\rho>0\)。由于 \(h\) 光滑且 \(K\) 紧，存在 \(C>0\)，使 \(h\le Cg_0\) 于 \(K\)。当 \(j\) 充分大时，从 \(v_j\) 到 \(v_\infty\) 的直线段含于 \(K\)，于是
	\[
		d_h(v_j,v_\infty)
		\le\sqrt C\,|v_j-v_\infty|_{g_0}\longrightarrow0.
	\]
	这就证明了完备性。
\end{proof}

\termidx[Cartan Hadamard dingli]{Cartan--Hadamard 定理}
\foreignidx{Cartan--Hadamard theorem}
\begin{theorem}[Cartan--Hadamard]\label{g3:thm:cartan-hadamard}
	设 \((M^n,g)\) 是连通、完备、光滑无边界黎曼流形，\(n\ge1\)，且截面曲率 \(K\le0\)，其中一维时曲率条件为空。则对每个 \(p\in M\)，
	\[
		\exp_p:\mathrm{T}_pM\longrightarrow M
	\]
	是光滑万有覆盖。若 \(M\) 还单连通，则 \(\exp_p\) 是全局微分同胚，故 \(M\) 微分同胚于 \(\mathbb R^n\)。
\end{theorem}
\begin{proof}
	\Cref{g3:prop:nonpositive-exponential-differential}使式\eqref{g3:eq:exponential-pullback-metric}定义了整个 \(\mathrm{T}_pM\) 上的黎曼度量 \(h\)，并且 \(h\ge g_0\)。前一引理证明 \((\mathrm{T}_pM,h)\) 完备。于是第五章的 \cref{g3:lem:complete-local-isometry-covering}适用于局部等距映射
	\[
		\exp_p:(\mathrm{T}_pM,h)\longrightarrow(M,g),
	\]
	得到它满射且为覆盖映射。该引理的关键条件是源流形完备：目标中的短径向测地线可以从每个原像点完整提升，从而构成互不相交的覆盖片。

	向量空间 \(\mathrm{T}_pM\) 单连通，因此这个覆盖是万有覆盖。若 \(M\) 也单连通，同一引理的最后一项说明连通覆盖只有一层，于是 \(\exp_p\) 是全局微分同胚。选择 \(\mathrm{T}_pM\) 的线性坐标即可得到它与 \(\mathbb R^n\) 的微分同胚。
\end{proof}

\begin{definition}\label{g3:def:hadamard-manifold}
	连通、完备、单连通、光滑无边界且截面曲率非正的黎曼流形称为 \term[Hadamard liuxing]{Hadamard 流形}[Hadamard manifold]。一维时包括通常的 Euclid 直线。
\end{definition}

拉回度量 \(h\) 通常并非常系数度量 \(g_0\)。定理给出的全局微分同胚允许用一个切空间为整个流形作坐标；它与 \(h\) 一起才是等距映射。非正曲率使这些坐标中的长度至少与 Euclid 长度一样大。

\begin{example}\label{g3:exm:hadamard-hypotheses}
	完备性和单连通性分别承担不同任务。Euclid 开单位球单连通且平坦，但测地线会在有限时间到达缺失的边界，指数映射不能定义在整个切空间上。平坦圆柱 \(S^1\times\mathbb R\) 完备且 \(K=0\)，它的指数映射则是从 \(\mathbb R^2\) 到圆柱的覆盖；相差圆周整倍长度的初速度有相同终点。这个例子说明，没有共轭点仍可有非空割迹。
\end{example}

\subsection{全局唯一测地线}
\label{g3:subsec:0901:03}

在 Hadamard 流形中，可把指数映射的逆记为
\[
	\log_p:M\longrightarrow\mathrm{T}_pM.
\]
它为每个终点记录唯一的测地初速度。为了以后对两个端点同时作变分，还需确认这条线随端点一起光滑变化。
\symidx{\(\log_p\)：指数映射在割前域上的逆，Hadamard 流形上为全局逆}

\begin{proposition}\label{g3:prop:hadamard-geodesic-uniqueness}
	设 \((M^n,g)\) 为 \(n\ge1\) 维 Hadamard 流形。记 \(\log_x=\exp_x^{-1}\)。对任意 \(x,y\in M\)，恰有一条定义于 \([0,1]\)、具有仿射参数且端点为 \(x,y\) 的测地线，公式为
	\begin{equation}\label{g3:eq:hadamard-geodesic}
		\gamma_{x,y}(t)=\exp_x(t\log_x y).
	\end{equation}
	它实现距离，且
	\[
		d(x,y)=|\log_x y|,\qquad
		d(\gamma_{x,y}(s),\gamma_{x,y}(t))
		=|s-t|\,d(x,y)\quad(s,t\in\mathbb R),
	\]
	其中测地线已按完备性延长到整个 \(\mathbb R\)。映射 \((x,y,t)\mapsto\gamma_{x,y}(t)\) 在 \(M\times M\times\mathbb R\) 上光滑。每个点的割迹均为空。
\end{proposition}
\begin{proof}
	由 Cartan--Hadamard 定理，\(\exp_x\) 是全局微分同胚，所以式\eqref{g3:eq:hadamard-geodesic}有定义。任意同端点、同参数区间的测地线都由某个初速度 \(v\) 决定，而其终点要求 \(\exp_x(v)=y\)。因此 \(v=\log_x y\)，得到唯一性。Hopf--Rinow 定理提供一条连接 \(x,y\) 的最短测地线，将它常速参数化到 \([0,1]\)，必与所写曲线一致。长度等于初速度的模，故 \(d(x,y)=|\log_x y|\)。

	这段论证可用于延长测地线上的任意两个点。因此它的每个有限子段都实现距离，且常速性给出 \(s,t\) 的距离等式。径向测地线始终最短，于是每个方向的割时都为无穷大，割迹为空。

	最后考虑
	\[
		\mathcal E:\mathrm{T}M\longrightarrow M\times M,\qquad
		(x,v)\longmapsto(x,\exp_xv).
	\]
	测地初值问题的光滑依赖性保证 \(\mathcal E\) 光滑；每个纤维上的指数映射双射保证它是集合双射。若 \(\mathrm d\mathcal E\) 的某个核向量投影到 \(\mathrm{T}_xM\)，第一分量说明该投影为零，所以这个核向量是竖直向量。第二分量此时就是 \((\mathrm d\exp_x)_v\) 作用于它，微分可逆性迫使它为零。因此 \(\mathcal E\) 的微分可逆，它的各局部光滑逆拼成全局光滑逆
	\[
		(x,y)\longmapsto(x,\log_x y).
	\]
	再与纤维内的标量乘法及全域指数映射复合，便得到 \(\gamma_{x,y}(t)\) 的联合光滑性。
\end{proof}

这里的唯一性包括所有具有指定仿射参数区间的测地线。若只把最短道路视为带任意速度的曲线，则还可作保持方向的重参数化。

\begin{proposition}\label{g3:prop:hadamard-exponential-expansion}
	设 \((M^n,g)\) 为 \(n\ge1\) 维 Hadamard 流形，\(p\in M\)。记 \(\log_p=\exp_p^{-1}\)，\(L\) 表示相应度量的曲线长度。对 \(\mathrm{T}_pM\) 中任意分片光滑曲线 \(c\)，有
	\[
		L_g(\exp_p\circ c)\ge L_{g_p}(c).
	\]
	而且对任意 \(v,w\in\mathrm{T}_pM\)，
	\begin{equation}\label{g3:eq:hadamard-distance-expansion}
		d(\exp_pv,\exp_pw)\ge|v-w|_{g_p}.
	\end{equation}
	等价地，\(\log_p\) 是从 \((M,d_g)\) 到 \((\mathrm{T}_pM,|\cdot|_{g_p})\) 的1-Lipschitz 映射。对径向点对，总有 \(d(p,\exp_pv)=|v|\)。
\end{proposition}
\begin{proof}
	沿曲线积分式\eqref{g3:eq:nonpositive-exponential-expansion}，得到长度不等式。为证明距离不等式，任取从 \(\exp_pv\) 到 \(\exp_pw\) 的分片光滑曲线 \(\alpha\)。全局光滑逆使 \(c=\log_p\circ\alpha\) 成为从 \(v\) 到 \(w\) 的曲线，故
	\[
		L_g(\alpha)\ge L_{g_p}(c)\ge|v-w|.
	\]
	对所有这样的 \(\alpha\) 取下确界，即得式\eqref{g3:eq:hadamard-distance-expansion}。最后一项已由径向线的最短性证明。
\end{proof}

此处全局逆保证了每一条竞争道路都能拉回到指定的两个初向量之间。仅考察切空间直线段的像，只能得到一条具体道路的长度，尚不足以控制目标流形中的距离下界。双曲空间中横向伸缩随半径增大；全局扩张估计也不意味着指数映射存在统一的 Lipschitz 上界。

\subsection{距离函数凸性和拓扑后果}
\label{g3:subsec:0901:04}

割迹为空以后，第七章的 Hessian 比较可以用于任意正半径。距离 \(r_p=d(p,\cdot)\) 仍在中心 \(p\) 不可微，但
\[
	F_p=\frac12r_p^2
\]
把这个尖点消去。其 Hessian 下界比通常的凸性更强。

\begin{proposition}\label{g3:prop:hadamard-squared-distance-convex}
	设 \((M^n,g)\) 为 \(n\ge1\) 维 Hadamard 流形，\(p\in M\)。记 \(r_p(q)=d(p,q)\)、\(\log_q=\exp_q^{-1}\)。函数 \(F_p(q)=\frac12d(p,q)^2\) 在整个 \(M\) 上光滑，满足
	\begin{equation}\label{g3:eq:hadamard-squared-distance-gradient}
		\operatorname{grad}F_p(q)=-\log_qp,\qquad
		\operatorname{Hess}F_p\ge g.
	\end{equation}
	在 \(M\setminus\{p\}\) 上，\(r_p\) 光滑且
	\begin{equation}\label{g3:eq:hadamard-distance-hessian}
		\operatorname{Hess}r_p
		\ge\frac1{r_p}\bigl(g-\mathrm dr_p\otimes\mathrm dr_p\bigr)
		\ge0.
	\end{equation}
	特别地，沿任意仿射参数测地线 \(\gamma\)，
	\[
		(F_p\circ\gamma)''\ge|\dot\gamma|^2.
	\]
\end{proposition}
\begin{proof}
	全局公式
	\[
		F_p(q)=\frac12|\log_pq|^2
	\]
	证明 \(F_p\) 光滑；范数在非零向量处光滑，故 \(r_p\) 在中心外光滑。由 \cref{g3:prop:hadamard-geodesic-uniqueness}，割迹为空，且径向线在所有半径上都最短。若 \(n\ge2\)，\cref{g3:thm:hessian-comparison}中的 \(k=0\) 情形因而适用于每个 \(q\ne p\)，给出式\eqref{g3:eq:hadamard-distance-hessian}的第一项。距离梯度为单位向量，所以右端非负。一维时，每段避开 \(p\) 的弧长坐标中 \(r_p\) 为仿射函数，且 \(g=\mathrm dr_p\otimes\mathrm dr_p\)，于是同式仍然成立。

	在中心外使用乘积法则，得到
	\[
		\operatorname{Hess}F_p
		=\mathrm dr_p\otimes\mathrm dr_p
			+r_p\operatorname{Hess}r_p\ge g.
	\]
	在 \(p\) 的法坐标中，\(F_p=\frac12\sum_i(x^i)^2\)，其微分在中心为零，二阶坐标导数为 \(\delta_{ij}\)，所以 \((\operatorname{Hess}F_p)_p=g_p\)。这也证明不等式包含中心。

	若 \(q\ne p\)，令 \(r=d(p,q)\)，并令 \(\eta:[0,r]\to M\) 为从 \(p\) 到 \(q\) 的单位速测地线。距离梯度公式给出
	\(\operatorname{grad}F_p(q)=r\dot\eta(r)\)。反向线段从 \(q\) 到 \(p\) 的时间一区间初速度为 \(-r\dot\eta(r)\)，因此
	\[
		\log_qp=-r\dot\eta(r).
	\]
	在 \(q=p\) 时，两边都为零，梯度公式仍成立。最后，沿仿射参数测地线有
	\[
		(F_p\circ\gamma)''
		=\operatorname{Hess}F_p(\dot\gamma,\dot\gamma)
			+\mathrm dF_p(\nabla_{\dot\gamma}\dot\gamma)
		=\operatorname{Hess}F_p(\dot\gamma,\dot\gamma),
	\]
	故得最后的估计。
\end{proof}

设 \(f\) 为 Hadamard 流形 \(M\) 上的函数。若对任意 \(x,y\in M\)，沿式\eqref{g3:eq:hadamard-geodesic}中的线段都有
\[
	f(\gamma_{x,y}(t))\le(1-t)f(x)+tf(y),\qquad 0\le t\le1.
\]
则称 \(f\) 为\term[ceditu hanshu]{测地凸函数}[geodesically convex function]。对光滑函数，Hessian 非负保证这一性质：沿线段的二阶导数非负，而区间上的凸函数不超过其两端值的线性插值。称集合为测地凸的，是指它包含任意两点之间的这条测地线段。

\begin{corollary}\label{g3:cor:hadamard-distance-balls-convex}
	设 \((M^n,g)\) 为 \(n\ge1\) 维 Hadamard 流形，\(p\in M\)。距离函数 \(r_p(q)=d(p,q)\) 与距离平方的一半 \(F_p(q)=d(p,q)^2/2\) 都是测地凸函数。对每个 \(R>0\)，开球 \(B(p,R)\) 和闭球 \(\overline B(p,R)\) 都测地凸，且闭球紧。
\end{corollary}
\begin{proof}
	上一命题已证明 \(F_p\) 的凸性。对 \(r_p\)，一条不经过 \(p\) 的测地线段上有 \((r_p\circ\gamma)''\ge0\)。若非恒定的仿射参数测地线在参数 \(t_0\) 经过 \(p\)，设其常速为 \(c>0\)。每个子段都最短，所以沿整条线有
	\[
		r_p(\gamma(t))=c|t-t_0|,
	\]
	这同样是凸函数。恒定线的结论直接成立，因此中心的不可微性没有破坏测地凸性。

	若两个端点的距离都小于 \(R\)，其线性插值也小于 \(R\)，凸性便使整个连接线段留在开球中。把严格不等式改为非严格不等式，得到闭球的结论。闭球紧性来自完备性和 Hopf--Rinow 定理。
\end{proof}

\begin{corollary}\label{g3:cor:hadamard-contractible}
	每个 \(n\ge1\) 维 Hadamard 流形 \(M\) 都可缩且非紧。对固定 \(p\in M\)，映射
	\[
		H:M\times[0,1]\longrightarrow M,\qquad
		H(q,s)=\exp_p\bigl((1-s)\log_pq\bigr)
	\]
	把 \(M\) 强形变收缩到 \(p\)。同一映射将每个以 \(p\) 为中心的开球或闭球收缩到 \(p\)。
\end{corollary}
\begin{proof}
	全局指数映射及其逆光滑，所以 \(H\) 光滑；它满足
	\(H(q,0)=q\)、\(H(q,1)=p\) 及 \(H(p,s)=p\)。此外，径向距离等式给出
	\[
		d(p,H(q,s))=(1-s)d(p,q),
	\]
	因而收缩始终保留各个开球或闭球。非紧性可由 \(M\) 微分同胚于 \(\mathbb R^n\) 得到，也可由任意单位径向线满足 \(d(p,\exp_p(tu))=t\) 看出。
\end{proof}

因此，正维、紧、单连通的流形不能携带截面曲率处处非正的黎曼度量。一般完备非正曲率流形的拓扑则集中在万有覆盖的不同覆盖片如何识别：其万有覆盖微分同胚于 Euclid 空间，而平坦环面已经给出紧的非单连通例子。距离平方的全局 Hessian 下界将在本章最后一节用来研究闭凸集的最近点及三角形比较。
% 纲要条目排为编号小节标题，并进入目录与书签；说明段仍保留于纲要注释中。

% !TeX root = ../../Book3.tex
% 纲要标题：### 9.2 常曲率的局部与整体分类
% 纲要位置：几何学计划纲要.md，第 2437 行。
% 正文状态：初稿；尚未正式审核。

% 纲要内容开始
% 
% * Jacobi 方程确定局部度量
% * 单连通完备模型
% * 离散自由等距群商
% * 球面、平坦与双曲空间形式
% 
% 纲要内容结束

\section{常曲率的局部与整体分类}
\label{g3:sec:0902}

常截面曲率把 Jacobi 方程中的曲率算子化为一个标量。因此，从一点发出的测地线怎样彼此分离，可以精确求解；这些分离量又决定极坐标中的角向度量。局部分类由此得到。整体分类还须处理两个不同的问题：正曲率下极坐标在对径点失效，非单连通情形下不同点会被覆盖变换识别。下面分别完成这两步。

\subsection{Jacobi 方程确定局部度量}
\label{g3:subsec:0902:01}

沿用\eqref{g3:eq:model-jacobi-functions}的 \(s_k\)，并记 \(b_k=\pi/\sqrt{k}\) 当 \(k>0\)，而 \(b_k=\infty\) 当 \(k\le0\)。记标准模型为
\begin{equation}\label{g3:eq:standard-space-form-models}
	M_k^n=
	\begin{cases}
		S^n(1/\sqrt{k}),&k>0,\\
		\mathbb R^n,&k=0,\\
		\bigl(\mathbb H^n,(-k)^{-1}g_{\mathbb H}\bigr),&k<0.
	\end{cases}
\end{equation}
这里 \(\mathbb H^n=\{(x,y):x\in\mathbb R^{n-1},\,y>0\}\)，且
\(g_{\mathbb H}=y^{-2}(|\mathrm dx|^2+\mathrm dy^2)\)。
\Cref{g3:exm:space-form-curvature,g3:exm:hyperbolic-warped-halfspace}已核算这些模型的曲率和完备性；后一个例子还用坐标估计证明了双曲度量的完备性。它们都连通，且当 \(n\ge2\) 时单连通：Euclid 空间和上半空间可缩，球面的情形是第一卷第5章第2节“高维球面的单连通性”命题。
% 跨卷依赖：g1:prop:sphere-simply-connected，已核读陈述及两开集 van Kampen 证明。
\symidx{\(M_k^n\)：单连通完备常曲率模型}

\begin{proposition}\label{g3:prop:space-form-local-model}
	取 \(k\in\mathbb R\)。记 \(s_k\) 为 \(s_k''+ks_k=0\)、\(s_k(0)=0\)、\(s_k'(0)=1\) 的解，\(b_k=\pi/\sqrt k\) 当 \(k>0\)，否则 \(b_k=\infty\)。
	设 \((M^n,g)\) 为光滑无边界黎曼流形，\(n\ge2\)，在 \(p\) 的某邻域中具有常截面曲率 \(k\)。在充分小的法球中，极坐标映射
	\[
		F_p(r,u)=\exp_p(ru),\qquad u\in S_pM,
	\]
	满足
	\begin{equation}\label{g3:eq:space-form-polar-metric}
		F_p^*g=\mathrm dr^2+s_k(r)^2g_{S_pM}.
	\end{equation}
	因此，若另一 \(n\) 维光滑无边界黎曼流形 \((N,h)\) 在 \(q\) 附近也具有常截面曲率 \(k\)，则任意线性等距
	\(I:\mathrm T_pM\to\mathrm T_qN\) 都确定一个局部等距
	\[
		\Phi=\exp_q\circ I\circ\exp_p^{-1},
		\qquad \Phi(p)=q,\quad \mathrm d\Phi_p=I.
	\]
	在充分小的连通邻域上，具有这些初始数据的局部等距唯一。
\end{proposition}
\begin{proof}
	取 \(\rho>0\)，使 \(\exp_p\) 在 \(B(0,\rho)\) 上为微分同胚，像位于所给常曲率邻域，且 \(\rho<b_k\)。固定单位向量 \(u\)，记 \(\gamma_u(r)=\exp_p(ru)\)、\(T=\dot\gamma_u\)，沿它的平行移动记为 \(P_r\)。对 \(w\in\mathrm T_uS_pM=u^\perp\)，取单位向量曲线 \(u(s)\)，使 \(u(0)=u\)、\(u'(0)=w\)。变分 \(\exp_p(ru(s))\) 的变分场满足
	\[
		J_w(0)=0,\qquad J_w'(0)=w,\qquad
		J_w''+R(J_w,T)T=0.
	\]
	常曲率给出 \(R(V,T)T=k(V-\langle V,T\rangle T)\)。由\cref{g3:exm:constant-curvature-jacobi}，或者直接在平行标架中解方程，得到
	\begin{equation}\label{g3:eq:space-form-angular-jacobi}
		J_w(r)=s_k(r)P_rw.
	\end{equation}
	因此角向内积为
	\(\langle J_w(r),J_z(r)\rangle=s_k(r)^2\langle w,z\rangle\)。
	径向速度为单位长度，又与每个 \(J_w\) 正交；后一事实也由 Gauss 引理给出。这就证明\eqref{g3:eq:space-form-polar-metric}。

	对 \(q\) 作同样计算，并进一步缩小 \(\rho\)。线性等距 \(I\) 保持单位切球面的度量，两边的极坐标公式完全相同，所以 \(\Phi\) 在去心邻域保持度量。它本身由光滑指数映射及其光滑局部逆组成，在中心也光滑，微分为 \(I\)，故中心处同样保持度量。

	最后，局部等距保持 Levi-Civita 联络，从而保持测地线及其仿射参数。任何具有所给初始数据的局部等距 \(\Psi\)，都满足
	\[
		\Psi(\exp_p v)=\exp_q(Iv)
	\]
	于充分小的切球上，因此与 \(\Phi\) 相同。这证明局部唯一性。
\end{proof}

极坐标公式的推导只需所考察径向测地线存在，并且沿线曲率为 \(k\)。所以当 \(M\) 完备且全局常曲率时，公式\eqref{g3:eq:space-form-angular-jacobi}在所有实参数上成立；当 \(0<r<b_k\) 时，径向与角向微分构成可逆的线性映射，\(\exp_p\) 因而局部可逆。不过，局部可逆尚未排除不同方向到达同一点。

\subsection{单连通完备模型}
\label{g3:subsec:0902:02}

当 \(k\le0\) 时，上一节的 Cartan--Hadamard 定理已经排除这种全局重复。正曲率情形恰好相反：所有角向 Jacobi 场在 \(b_k\) 处同时归零。我们需要证明，这种退化与标准球面南极的退化完全相同。

\begin{lemma}\label{g3:lem:round-unit-sphere-isometry}
	设 \(E,F\) 为维数相同的实内积空间，\(\dim E=\dim F=n\ge2\)。若光滑映射 \(v:S(E)\to S(F)\) 是单位球面之间的局部黎曼等距，则存在唯一线性等距 \(A:E\to F\)，使 \(v(u)=Au\) 对所有 \(u\in S(E)\) 成立。
\end{lemma}
\begin{proof}
	固定 \(u_0\in S(E)\)。定义
	\[
		Au_0=v(u_0),\qquad
		A|_{u_0^\perp}=\mathrm dv_{u_0}.
	\]
	右侧将 \(u_0^\perp\) 等距映到 \(v(u_0)^\perp\)，因此 \(A\) 是线性等距。对任意单位向量 \(w\perp u_0\)，曲线
	\(c(t)=\cos t\,u_0+\sin t\,w\)
	是单位球面上的单位速大圆。局部等距保持测地线，所以 \(v\circ c\) 是初值
	\((v(u_0),\mathrm dv_{u_0}w)\)
	的球面测地线。球面测地线的显式解及初值唯一性给出
	\[
		v(c(t))=\cos t\,v(u_0)+\sin t\,\mathrm dv_{u_0}w=Ac(t).
	\]
	每个不同于 \(-u_0\) 的球面点都能如此表示；在 \(-u_0\) 处由连续性仍成立。因此 \(v=A|_{S(E)}\)。单位向量张成 \(E\)，又保证线性延拓唯一。
\end{proof}

\begin{theorem}\label{g3:thm:simply-connected-space-form}
	设 \((M^n,g)\) 为连通、单连通、完备、光滑无边界黎曼流形，\(n\ge2\)，且截面曲率处处等于常数 \(k\)。则 \(M\) 等距于\eqref{g3:eq:standard-space-form-models}中的 \(M_k^n\)。更准确地，给定 \(o\in M_k^n\)、\(p\in M\) 及线性等距 \(I:\mathrm T_oM_k^n\to\mathrm T_pM\)，存在唯一整体等距 \(F:M_k^n\to M\)，满足 \(F(o)=p\)、\(\mathrm dF_o=I\)。
\end{theorem}
\begin{proof}
	先设 \(k\le0\)。\Cref{g3:thm:cartan-hadamard}分别应用于 \(M\) 和 \(M_k^n\)，说明两个基点的指数映射都是全局微分同胚。因此
	\[
		F=\exp_p\circ I\circ(\exp_o)^{-1}
	\]
	全局有定义。两边在所有 \(r>0\) 上都有\eqref{g3:eq:space-form-polar-metric}，而 \(I\) 保持角向单位球度量，所以 \(F\) 处处保持度量；在中心的光滑性由正常坐标直接给出。这就得到所需整体等距。

	以下设 \(k>0\)，记 \(a=1/\sqrt{k}\)、\(L=\pi a\)，模型为 \(S^n(a)\)。取 \(o\) 的对径点 \(o_-\)。在模型去掉 \(o_-\) 后，每点唯一写成 \(\exp_o(ru)\)，其中 \(0\le r<L\)、\(|u|=1\)。定义
	\[
		F(\exp_o(ru))=\gamma_{Iu}(r),
		\qquad \gamma_\xi(r)=\exp_p(r\xi).
	\]
	完备性使右侧的测地线全时存在。局部极坐标公式说明 \(F\) 在这一去点球面上为局部等距，包括北极 \(o\)。

	先找出南极应当映到哪里。映射
	\[
		Q:S_pM\to M,\qquad Q(\xi)=\exp_p(L\xi)
	\]
	光滑，且对 \(\eta\perp\xi\)，其微分为
	\(\mathrm dQ_\xi(\eta)=J_\eta(L)=s_k(L)P_L\eta=0\)。
	因为 \(n\ge2\)，单位球面 \(S_pM\) 道路连通；沿球面中的每条光滑道路，\(Q\) 的导数恒为零，所以 \(Q\) 是常映射。记其值为 \(q\)。这说明所有从 \(p\) 出发的单位速径向线在时刻 \(L\) 都到达同一点 \(q\)。

	还须检查在 \(q\) 附近怎样离开这个公共终点。定义反向末速度映射
	\[
		v:S_pM\to S_qM,\qquad
		v(\xi)=-\dot\gamma_\xi(L).
	\]
	它是光滑的，且值为单位向量。对初方向变分 \(\xi(s)\)，因为终点恒为 \(q\)，对 \(v\) 的普通纤维微分就是相应的协变导数。无挠性给出混合导数交换，从而
	\[
		\mathrm dv_\xi(\eta)
		=-J_\eta'(L)
		=-c_k(L)P_L\eta
		=P_L\eta.
	\]
	这里 \(c_k(L)=\cos\pi=-1\)。该式保持内积，且其像垂直于 \(v(\xi)\)，所以 \(v\) 是单位球面之间的局部等距。由上一引理，存在一个线性等距
	\(A:\mathrm T_pM\to\mathrm T_qM\)，使 \(v(\xi)=A\xi\)。

	在标准球面上同样记
	\[
		v_0(u)=-\left.\frac{\mathrm d}{\mathrm dr}\exp_o(ru)\right|_{r=L}
		\in S_{o_-}S^n(a).
	\]
	大圆公式表明 \(v_0\) 是某个线性等距
	\(A_0:\mathrm T_oS^n(a)\to\mathrm T_{o_-}S^n(a)\)
	在单位球面上的限制。对充分小的 \(s>0\)，反向测地线的初值唯一性给出
	\[
		\exp_o((L-s)u)=\exp_{o_-}(sA_0u),
		\qquad
		\gamma_{Iu}(L-s)=\exp_q(sAIu).
	\]
	于是原来的 \(F\) 在 \(o_-\) 的去心正常邻域内恰等于
	\begin{equation}\label{g3:eq:space-form-antipodal-extension}
		\exp_q\circ(AIA_0^{-1})\circ(\exp_{o_-})^{-1}.
	\end{equation}
	右侧在中心也光滑，且由\cref{g3:prop:space-form-local-model}是局部等距。令 \(F(o_-)=q\)，便把 \(F\) 延拓为整个 \(S^n(a)\) 上的光滑局部等距。这里不仅证明了径向像收敛，还确定了中心处的微分 \(AIA_0^{-1}\)。

	圆球面完备，而 \(M\) 连通、单连通，故\cref{g3:lem:complete-local-isometry-covering}使这个全域局部等距成为整体等距同构。最后，对任意符号的 \(k\)，具有初始数据 \((p,I)\) 的等距都满足
	\(F(\exp_o z)=\exp_p(Iz)\)。
	模型完备，Hopf--Rinow 定理保证 \(\exp_o\) 满射，所以这条等式唯一确定 \(F\)。
\end{proof}

在正曲率证明中，完备性首先保证测地线能走到共同终点，并在该点反向延伸；南极附近的光滑性则由末速度映射保证。这个区分很重要：仅知道所有角向微分在某个时刻归零，还不足以直接断言已经得到了一个光滑球面。

\subsection{离散自由等距群商}
\label{g3:subsec:0902:03}

单连通模型确定以后，一般空间形式的差异由哪些模型点被识别来记录。先明确群作用的条件。在第一卷第5章第5节中，作用“适当不连续”采用局部覆盖片的约定：每个点有一个开邻域 \(U\)，满足
\[
	\gamma U\cap U=\varnothing\qquad(\gamma\ne e).
\]
它已经包含自由性。对于完备黎曼流形的等距作用，这一条件等价于“作用自由，且对每个紧集 \(C\)，满足 \(\gamma C\cap C\ne\varnothing\) 的群元素只有有限个”。下面同时验证这两种实际使用的形式。

\begin{lemma}\label{g3:lem:space-form-quotient-action}
	设 \((X,g)\) 为连通、完备、光滑无边界黎曼流形，\(\Gamma\) 为其等距变换群的一个子群。则以下条件等价：
	\begin{equivalences}
		\item 每个 \(x\in X\) 都有开邻域 \(U\)，其不同 \(\Gamma\) 平移两两不交。
		\item \(\Gamma\) 自由作用，且每个紧集 \(C\subset X\) 都满足
		\[
			\#\{\gamma\in\Gamma:\gamma C\cap C\ne\varnothing\}<\infty.
		\]
	\end{equivalences}
	在这些条件下，\(\Gamma\) 在等距群的紧开拓扑中是离散子群，商空间 \(X/\Gamma\) 有唯一使投影局部等距的光滑黎曼结构，而且此商连通、完备且无边界。
\end{lemma}
\begin{proof}
	设第一条成立。若 \(\gamma x=x\)，局部不交性立即给出 \(\gamma=e\)，所以作用自由。固定 \(x_0\)，在一个这样的邻域内取 \(B(x_0,\varepsilon)\)。其不同平移不交；由测地线中点可知，不同轨道点之间的距离至少为 \(2\varepsilon\)。设 \(C\) 包含于 \(\overline B(x_0,R)\)。若 \(\gamma C\cap C\ne\varnothing\)，选 \(z\in C\) 使 \(\gamma z\in C\)，则
	\[
		d(x_0,\gamma x_0)
		\le d(x_0,\gamma z)+d(\gamma z,\gamma x_0)
		\le2R.
	\]
	闭球 \(\overline B(x_0,2R)\) 紧，其中的 \(2\varepsilon\) 分离点集有限；自由性使轨道点与群元素一一对应。这就得到第二条。

	反过来，固定 \(x\)，取紧闭球 \(C=\overline B(x,1)\)。只有有限个非恒等元素可能满足 \(\gamma C\cap C\ne\varnothing\)。由自由性，这些元素的 \(d(x,\gamma x)\) 都正。取 \(\delta<1\)，使 \(2\delta\) 小于这些有限个正数的最小值；若集合为空，只需取 \(\delta<1\)。那么 \(U=B(x,\delta)\) 与所有非恒等平移不交：若有交点，先知 \(\gamma C\cap C\ne\varnothing\)，再由三角不等式得到 \(d(x,\gamma x)<2\delta\)，矛盾。

	局部不交性还给出一个仅含恒等元的群邻域：选上述 \(x,\delta\)，则
	\(\{\gamma:d(x,\gamma x)<2\delta\}\)
	在紧开拓扑中开，并与 \(\Gamma\) 只交于 \(e\)。故 \(\Gamma\) 离散。第一卷第5章第5节“作用的商是覆盖”命题说明，商投影 \(\pi:X\to X/\Gamma\) 为覆盖，其覆盖片就是这些 \(U\) 及其平移。
	% 跨卷依赖：g1:prop:action-quotient-cover，已核读局部片构造及覆盖变换群识别。

	还须核对商的流形分离条件。若 \(x,y\) 不在同一轨道，取紧集 \(C=\overline B(x,1)\cup\{y\}\)。满足 \(d(x,\gamma y)\le1\) 的 \(\gamma\) 只有有限个，且每个距离都正，所以
	\(\inf_{\gamma\in\Gamma}d(x,\gamma y)>0\)。
	取半径小于这个下界三分之一的两球，其投影为不交的开邻域，故商 Hausdorff。投影是开映射，\(X\) 的一个可数拓扑基的像便给出商的可数拓扑基。将每张片的光滑坐标搬到商上，换片映射局部等于某个 \(\gamma\)，因而光滑；它们又是等距，故各片的度量相容。这构造出唯一使 \(\pi\) 局部等距的光滑度量。连通性和无边界性直接由投影及局部片得到。最后，源完备且 \(\pi\) 是黎曼覆盖，\cref{g3:prop:riemannian-cover-complete}给出商完备。
\end{proof}

\begin{theorem}\label{g3:thm:space-form-quotient}
	设 \(n\ge2\)、\(k\in\mathbb R\)。连通、完备、光滑无边界的 \(n\) 维常截面曲率 \(k\) 黎曼流形，恰好是
	\[
		M_k^n/\Gamma,
	\]
	式中 \(M_k^n\) 为式\eqref{g3:eq:standard-space-form-models}的单连通完备标准模型，\(\Gamma\subset\operatorname{Isom}(M_k^n)\) 自由且适当不连续地作用；这样的 \(\Gamma\) 自动离散，并满足\cref{g3:lem:space-form-quotient-action}的紧集有限性条件。商的万有覆盖为 \(M_k^n\)，基本群同构于 \(\Gamma\)。
\end{theorem}
\begin{proof}
	给定空间形式 \(M\)，由\cref{g3:prop:riemannian-universal-cover}取带拉回度量的万有覆盖 \(\widetilde M\)。它连通、单连通、完备，且局部等距保持曲率，所以截面曲率仍为 \(k\)。由\cref{g3:thm:simply-connected-space-form}，存在等距 \(F:M_k^n\to\widetilde M\)。

	将 \(\widetilde M\to M\) 的覆盖变换群经 \(F\) 搬到模型上，得到等距子群 \(\Gamma\)。均匀覆盖邻域的一张连通片与其不同覆盖变换的像不交：若两片相交，它们相同，而片上投影单射，迫使该变换在片上恒等；覆盖变换由一点决定，所以它全局恒等。因此 \(\Gamma\) 满足上一引理的局部条件。万有覆盖的变换群在每个纤维上自由传递，故投影的纤维恰为这些轨道。于是 \(\pi\circ F\) 诱导从 \(M_k^n/\Gamma\) 到 \(M\) 的双射局部等距，因而是整体等距。其基本群与 \(\Gamma\) 的同构同样来自万有覆盖。

	反过来，给定这样的 \(\Gamma\)，上一引理已经构造出连通、完备、光滑无边界的商度量。商投影局部等距，所以商具有常曲率 \(k\)。模型单连通，因此该覆盖是万有覆盖；第一卷的商覆盖命题还说明它的全部覆盖变换恰是给定的 \(\Gamma\)，由此得到基本群。
\end{proof}

群的抽象同构类型不能单独决定商度量；同一个群可以通过不同尺度、不同平移向量作用。等距分类中需要保留它作为模型等距群子群的位置。

\begin{proposition}\label{g3:prop:space-form-isometry-conjugacy}
	设 \(n\ge2\)、\(k\in\mathbb R\)，以 \(M_k^n\) 表示式\eqref{g3:eq:standard-space-form-models}的单连通完备标准模型，且 \(\Gamma_1,\Gamma_2\subset\operatorname{Isom}(M_k^n)\) 都自由、适当不连续地作用。则两个黎曼商 \(M_k^n/\Gamma_1\) 与 \(M_k^n/\Gamma_2\) 等距，当且仅当存在 \(A\in\operatorname{Isom}(M_k^n)\)，使
	\[
		A\Gamma_1A^{-1}=\Gamma_2.
	\]
\end{proposition}
\begin{proof}
	若有这样的 \(A\)，定义 \([x]_{\Gamma_1}\mapsto[Ax]_{\Gamma_2}\)。共轭关系保证它良定义，\(A^{-1}\) 给出逆映射，而在覆盖片上它就是 \(A\)，故为等距。

	反过来，设 \(f:M_k^n/\Gamma_1\to M_k^n/\Gamma_2\) 为黎曼等距，商投影记为 \(\pi_1,\pi_2\)。固定 \(x_0\in M_k^n\)，再取 \(y_0\in\pi_2^{-1}(f(\pi_1(x_0)))\)。第一卷第5章第3节“映射的提升判据”应用于 \(f\circ\pi_1\) 和覆盖 \(\pi_2\)：定义域单连通，子群包含条件自动满足，故存在唯一提升 \(A\)，满足
	\[
		\pi_2\circ A=f\circ\pi_1,\qquad A(x_0)=y_0.
	\]
	% 跨卷依赖：g1:thm:lifting-criterion，已核读存在性、唯一性及覆盖片内的局部表达。
	在覆盖片上，\(A\) 是两个局部覆盖逆及 \(f\) 的复合，所以光滑且局部等距。用同一判据提升 \(f^{-1}\circ\pi_2\)，指定 \(B(y_0)=x_0\)。于是 \(BA\) 与恒等映射都是 \(\pi_1\) 的提升，且在 \(x_0\) 处相同，故 \(BA=\mathrm{id}\)；同理 \(AB=\mathrm{id}\)。因此 \(A\) 是整体等距。

	对每个 \(\gamma\in\Gamma_1\)，有
	\[
		\pi_2\circ A\gamma A^{-1}
		=f\circ\pi_1\circ\gamma A^{-1}
		=f\circ\pi_1\circ A^{-1}=\pi_2.
	\]
	所以 \(A\gamma A^{-1}\) 是 \(\pi_2\) 的覆盖变换，属于 \(\Gamma_2\)。将同样论证用于 \(A^{-1}\)，得到反向包含，因而两群共轭。
\end{proof}

\subsection{球面、平坦与双曲空间形式}
\label{g3:subsec:0902:04}

分类定理把局部度量确定以后剩余的自由度集中到离散群作用上。下面的模型既给出不同拓扑，也显示正曲率与非正曲率的商有不同的整体规模。

\begin{example}\label{g3:exm:spherical-space-form-quotients}
	当 \(k>0\) 时，模型 \(S^n(a)\)、\(a=1/\sqrt{k}\) 紧。将紧集有限性用于整个球面，便知允许的 \(\Gamma\) 必有限。因此每个正曲率空间形式都是紧流形，且基本群有限，这与 Bonnet--Myers 定理一致。

	最简单的非平凡作用由对径映射 \(x\mapsto-x\) 生成。它是等距且没有不动点，有限性保证紧集条件，所以
	\[
		S^n(a)/\{I,-I\}=\mathbb {RP}^n
	\]
	带有常曲率 \(a^{-2}\) 的完备度量。由球面距离公式及覆盖距离公式，对两条实直线的单位代表 \(u,v\)，有
	\[
		d_{\mathbb {RP}^n}([au],[av])
		=\min\{a\arccos\langle u,v\rangle,\,
		a\arccos(-\langle u,v\rangle)\}
		=a\arccos|\langle u,v\rangle|.
	\]
	故其直径为 \(\pi a/2\)，小于覆盖球面的 \(\pi a\)。球面模型与商空间曲率相同，最远点的距离却已改变。
\end{example}

\begin{example}\label{g3:exm:flat-torus-klein-spaceforms}
	在 \(\mathbb R^n\) 中取线性无关向量 \(v_1,\ldots,v_n\)，它们的整数线性组合构成满秩格
	\(\Lambda=\sum_i\mathbb Zv_i\)。
	平移作用自由。给定紧集 \(C\)，若 \((C+\lambda)\cap C\ne\varnothing\)，则 \(\lambda\in C-C\)，而有界集只含有限个格点，因此作用适当不连续。商
	\[
		T_\Lambda^n=\mathbb R^n/\Lambda
	\]
	是完备平坦环面。闭平行多面体
	\(\{\sum_it_iv_i:0\le t_i\le1\}\)
	投影覆盖整个商，所以该商紧。

	自由等距作用还可以包含线性反射部分。在 \(\mathbb R^2\) 中取 \(a,b>0\)，令
	\[
		A(x,y)=(x+a,y),\qquad
		B(x,y)=(-x,y+b).
	\]
	有 \(BAB^{-1}=A^{-1}\)，所以每个群元素都可写成 \(A^mB^j\)，其作用为
	\[
		A^mB^j(x,y)=((-1)^jx+ma,y+jb).
	\]
	这些表达互不相同。若有不动点，先由第二坐标得到 \(j=0\)，再由第一坐标得到 \(m=0\)，故作用自由。若某个固定紧集与其像相交，第二坐标限制 \(j\) 只能取有限个值，第一坐标再限制 \(m\) 只能取有限个值，故作用适当不连续。所得商是完备平坦 Klein 瓶：可先用 \(B\) 把 \(y\) 移入 \([0,b]\)，再用 \(A\) 把 \(x\) 移入 \([0,a]\)，由矩形边界识别得到其通常模型。它与平坦环面具有相同的局部度量，但 \(B\) 反转定向，而格平移均保持定向。
\end{example}

\begin{example}\label{g3:exm:hyperbolic-cyclic-spaceform}
	在曲率 \(-1\) 的上半平面
	\[
		\mathbb H^2=\{(x,y):y>0\},\qquad
		g=\frac{\mathrm dx^2+\mathrm dy^2}{y^2}
	\]
	中，固定 \(\ell>0\)，令
	\(D_\ell(x,y)=(e^\ell x,e^\ell y)\)。
	分子、分母同时乘以 \(e^{2\ell}\)，所以 \(D_\ell\) 等距。它的非零次幂没有不动点。若紧集 \(C\) 的 \(y\) 坐标介于 \(c>0\) 与 \(C_0<\infty\) 之间，则
	\(D_\ell^jC\cap C\ne\varnothing\)
	迫使
	\[
		\frac c{C_0}\le e^{j\ell}\le\frac{C_0}c,
	\]
	所以只有有限个整数 \(j\) 可能出现。分类定理因此给出完备双曲曲面
	\(\mathbb H^2/\langle D_\ell\rangle\)。

	全局坐标 \(s=x/y\)、\(t=\log y\) 将 \(D_\ell\) 写成
	\((s,t)\mapsto(s,t+\ell)\)，所以商同胚于
	\(\mathbb R\times(\mathbb R/\ell\mathbb Z)\)，是一个圆柱。竖直曲线
	\(\gamma(t)=(0,e^t)\) 是单位速测地线，其长度最短性也可直接看出：任意曲线的速度满足
	\(|(\log y)'|\le|\dot\gamma|_g\)，连接两个竖直点的曲线长度至少为它们的 \(\log y\) 差，而竖直线达到下界。该轴在商中成为长度为 \(\ell\) 的闭测地线。

	任意代表 \(j\) 次生成元的闭曲线提升后，从某个 \((x,y)\) 到达
	\(D_\ell^j(x,y)\)，长度至少为 \(|j|\ell\)。因此这个商中最短非平凡闭曲线的长度为 \(\ell\)。改变 \(\ell\) 保持基本群同构于 \(\mathbb Z\)，却改变这一等距不变量，得到彼此不等距的完备双曲空间形式。
\end{example}

这些例子中的自由性保证商仍是无边界光滑流形；若去掉自由性，反射等距可能产生边界，旋转等距可能产生锥点。空间形式分类所讨论的是光滑黎曼流形，因此群作用的条件属于分类结论的一部分。
% 纲要条目排为编号小节标题，并进入目录与书签；说明段仍保留于纲要注释中。

% !TeX root = ../../Book3.tex
% 纲要标题：### 9.3* Toponogov 三角比较
% 纲要位置：几何学计划纲要.md，第 2444 行。
% 正文状态：初稿；尚未正式审核。

% 纲要内容开始
% 
% * 铰链比较与模型角
% * 局部比较及拼接引理
% * 全局三角比较的证明
% * 正曲率模型中的周长限制
% * 等号和刚性情形
% 
% 纲要内容结束

\OptionalSection{Toponogov 三角比较}{g3:sec:0903}

固定两条边及其夹角时，正曲率会使两条径向测地线的端点更接近；固定三条边时，同一现象则表现为角更大、顶点到对边的距离更大。第七章已经在距离光滑处给出了相应的 Hessian 比较。本节要处理的是长三角形：对边可能经过割迹，距离不再可微，局部公式不能直接沿整条边积分。

这里采用一维支撑函数的方法。先把距离换成一个适合模型方程的函数，再证明局部上支撑不等式能够在整段区间上拼合。这个证明允许对边与割迹任意相交，不要求把它划成有限个光滑段。所得结果是 Toponogov 三角比较。\footnote{\OriginalName{Виктор Андреевич Топоногов}（Victor Andreevich Toponogov，1930--2004），俄罗斯数学家，研究整体黎曼几何，建立了曲率下界的三角比较及非负曲率分裂结果。}

\subsection{铰链比较与模型角}
\label{g3:subsec:0903:01}

以 \(M_k^2\) 表示常截面曲率 \(k\) 的单连通完备二维模型：\(k>0\) 时为半径 \(1/\sqrt{k}\) 的球面，\(k=0\) 时为平面，\(k<0\) 时为双曲平面。继续使用第七章的 \(s_k,c_k\) 与 \(b_k\)，特别是
\[
 s_k'=c_k,\qquad c_k'=-ks_k,\qquad c_k^2+ks_k^2=1,
 \qquad b_k=\begin{cases}\pi/\sqrt{k},&k>0,\\ \infty,&k\le0.\end{cases}
\]
为了同时处理三个模型，引入函数
\begin{equation}\label{g3:eq:modified-model-distance}
 m_k(r)=\int_0^r s_k(t)\,\mathrm dt
 =\begin{cases}
 (1-c_k(r))/k,&k\ne0,\\ r^2/2,&k=0.
 \end{cases}
 \qquad m_k'=s_k,\quad m_k''=c_k=1-km_k.
\end{equation}
它在 \([0,b_k)\) 上严格递增，在零处有展开 \(m_k(r)=r^2/2+O(r^4)\)。将它延拓到负参数时，\(m_k\) 是偶函数；下文的 \(m_k(a-b)\) 因而也有意义。
\symidx{\(m_k(r)\)：模型修改距离函数，等于 \(\int_0^r s_k(t)\,\mathrm dt\)}

\begin{definition}\label{g3:def:comparison-triangle}
 设 \(p,q,r\) 为黎曼流形中三个不同的点，并分别选定连接它们的最短测地线段。这三个线段组成一个\term[ce di sanjiao xing]{测地三角形}[geodesic triangle]。记
 \[
 a=d(q,r),\qquad b=d(p,r),\qquad c=d(p,q),\qquad P=a+b+c.
 \]
 若 \(P<2b_k\)，则在 \(M_k^2\) 中具有相同三边长的三角形称为它的 \(k\)-\term[bijiao sanjiao xing]{比较三角形}[comparison triangle]，顶点记为 \(\bar p,\bar q,\bar r\)。其在 \(\bar p\) 处的角记为 \(\widetilde\angle_k(qpr)\)，称为相应的\term[moxing jiao]{模型角}[comparison angle]。允许三角不等式取等号，此时模型三角形退化为线段。

 两条从同一点出发的最短测地线段及其初始夹角组成一个\term[jiaolian]{铰链}[hinge]。同腿长、同夹角的模型铰链称为其比较铰链。
\end{definition}

比较三角形的存在与唯一性可以直接从余弦公式看出。在模型中，两条长度为 \(b,c\) 的腿形成角 \(\alpha\)，其端点距离 \(a\) 满足
\begin{equation}\label{g3:eq:model-hinge-cosine}
 m_k(a)=m_k(b-c)+s_k(b)s_k(c)(1-\cos\alpha).
\end{equation}
当 \(k=0\) 时，这就是 \(a^2=b^2+c^2-2bc\cos\alpha\)。当 \(k\ne0\) 时，它等价于
\[
 c_k(a)=c_k(b)c_k(c)+k s_k(b)s_k(c)\cos\alpha.
\]
在球面嵌入模型中取两条子午线，端点的 Euclid 内积给出这个等式；在双曲面模型中用 Lorentz 内积作同一计算，空间分量的系数变为 \(k<0\)，得到双曲余弦公式。这也给出统一公式\eqref{g3:eq:model-hinge-cosine}。

\ProseParagraphBreak

若 \(0<b,c<b_k\)，模型第三边随 \(\alpha\in[0,\pi]\) 连续增加。\(k\le0\) 时，其范围是 \([|b-c|,b+c]\)；\(k>0\) 时，球面的最短距离使范围成为
\[
 [\,|b-c|,\ \min\{b+c,2b_k-b-c\}\,].
\]
三角不等式与 \(P<2b_k\) 正好保证给定的 \(a\) 位于这个范围内。因此比较三角形存在，在模型等距意义下唯一；非退化时三个模型角都在 \((0,\pi)\) 内。

\ProseParagraphBreak

需要区分两种固定数据。固定三边长时，曲率下界将给出“实际角不小于模型角”；固定两腿与夹角时，它将给出“实际第三边不大于模型第三边”。两句话由式\eqref{g3:eq:model-hinge-cosine}的单调性相互对应。以下先证明三边形式，因为它自然成为一个端点值已知的一维比较问题。

\subsection{局部比较及拼接引理}
\label{g3:subsec:0903:02}

设 \(r_p=d(p,\cdot)\) 在某点光滑，并且 \(0<r_p<b_k\)。若 \(K\ge k\)，\cref{g3:thm:hessian-comparison}给出
\[
 \operatorname{Hess}r_p\le\frac{c_k(r_p)}{s_k(r_p)}
 (g-\mathrm dr_p\otimes\mathrm dr_p).
\]
对复合函数 \(m_k\circ r_p\) 求 Hessian，径向项和横向项恰好合在一起：
\begin{equation}\label{g3:eq:modified-distance-hessian-smooth}
 \begin{split}
 \operatorname{Hess}(m_k\circ r_p)
 &=c_k(r_p)\,\mathrm dr_p\otimes\mathrm dr_p
   +s_k(r_p)\operatorname{Hess}r_p\\
 &\le c_k(r_p)g=(1-km_k(r_p))g.
 \end{split}
\end{equation}
因而沿单位速测地线 \(\sigma\)，函数 \(h=m_k\circ r_p\circ\sigma\) 满足 \(h''+kh\le1\)。模型中的对应函数满足等式。

\ProseParagraphBreak

在割迹上使用这个不等式，需要给“\(h''\)”一个不要求可微的准确含义。若连续函数 \(h\) 在内点 \(t_0\) 满足：对每个 \(\varepsilon>0\)，都有 \(t_0\) 邻域上的 \(C^2\) 函数 \(H\)，使
\begin{equation}\label{g3:eq:one-dimensional-upper-support}
 H\ge h,\qquad H(t_0)=h(t_0),\qquad
 H''(t_0)+kh(t_0)\le1+\varepsilon,
\end{equation}
就说 \(h''+kh\le1\) 在 \(t_0\) 以\term[shang zhicheng]{上支撑}[upper support]的意义成立。支撑只须在这个点的一个邻域内定义；邻域及支撑函数可以随点和 \(\varepsilon\) 改变。

\begin{lemma}[一维支撑拼接]\label{g3:lem:one-dimensional-support-gluing}
 设 \(k\in\mathbb R\)、\(0<L<b_k\)，\(h\in C([0,L])\) 在每个内点以上支撑意义满足 \(h''+kh\le1\)。设 \(\bar h\in C^2([0,L])\) 满足
 \[
 \bar h''+k\bar h=1,\qquad \bar h(0)=h(0),\quad \bar h(L)=h(L).
 \]
 则 \(h\ge\bar h\) 于整个 \([0,L]\)。更一般地，端点只满足 \(h\ge\bar h\) 时，结论仍成立。

 若 \(h\) 连续且分段 \(C^2\)，各段上满足相同微分不等式，且每个拼接点的单侧导数满足
 \begin{equation}\label{g3:eq:gluing-slope-condition}
 h'_+(t_j)\le h'_-(t_j),
 \end{equation}
 则它满足上述支撑条件，因而也满足全区间比较。
\end{lemma}
\begin{proof}
 先证明支撑条件足以给出全局结论。令 \(w=h-\bar h\)。它的上支撑满足 \(W''(t_0)+kw(t_0)\le\varepsilon\)。选择常数 \(\lambda>0\)，使
 \[
 \max\{k,0\}<\lambda^2<(\pi/L)^2,
 \qquad \psi(t)=\cos\bigl(\lambda(t-L/2)\bigr)>0.
 \]
 这样的 \(\lambda\) 存在，正是因为 \(L<b_k\)。若 \(w\) 有负值，则 \(w/\psi\) 在某内点 \(t_0\) 取负最小值 \(A<0\)：端点值非负，故最小点不能在端点。于是
 \(w-A\psi\ge0\)，且在 \(t_0\) 取零。

 取 \(w\) 在 \(t_0\) 的一个上支撑 \(W\)。由于 \(W\ge w\)，函数 \(W-A\psi\) 也在 \(t_0\) 取局部最小值，所以
 \[
 0\le W''(t_0)-A\psi''(t_0)
 \le\varepsilon+A(\lambda^2-k)\psi(t_0).
 \]
 右侧的第二项严格为负，选取足够小的 \(\varepsilon\) 即得矛盾。这证明 \(w\ge0\)。它同时说明：每一点的局部支撑，即使来自不同的构造，也足以比较整条区间；无须先选择一个共同的光滑分解。

 最后核对分段拼接条件。若某点满足严格跳跃 \(h'_+<h'_-\)，选择斜率 \(v\) 严格夹在两者之间。对任意固定实数 \(C\)，二次函数
 \[
 H(t_j+s)=h(t_j)+vs+\tfrac12Cs^2
 \]
 在充分小的双侧邻域内高于 \(h\)：右侧的正线性差来自 \(v>h'_+\)，左侧来自 \(v<h'_-\)。因此可取 \(C=1-kh(t_j)\)。若两侧导数相等，则两侧的二阶导数都不大于 \(1-kh(t_j)\)；Taylor 公式说明取
 \(C=1-kh(t_j)+\varepsilon\) 得到局部上支撑。这里分段 \(C^2\) 指各段连同端点具有单侧二阶连续延拓。在其余点直接取 \(H=h\)。故拼接后的函数满足式\eqref{g3:eq:one-dimensional-upper-support}。
\end{proof}

\(k=0\) 时，这个引理说：二阶导数不大于一的函数位于具有同样端点的模型抛物线之上。式\eqref{g3:eq:gluing-slope-condition}允许斜率向下跳；向上跳则会产生一个不能由光滑函数从上方接触的尖点。距离在割迹处切换最短分支，恰好适合前一种支撑结构。

\begin{proposition}[割前的局部三角比较]\label{g3:prop:local-toponogov}
 设 \((M^n,g)\) 连通、完备、光滑且无边界，\(n\ge2\)、\(K\ge k\)。取最短测地三角形 \(pqr\)，周长小于 \(2b_k\)，令 \(\sigma:[0,a]\to M\) 为从 \(q\) 到 \(r\) 的单位速边。若 \(\sigma([0,a])\) 不交 \(p\) 的割迹，则对相应的模型边 \(\bar\sigma\)，有
 \[
 d(p,\sigma(t))\ge d_k(\bar p,\bar\sigma(t))\qquad(0\le t\le a).
 \]
 若整个三角形处于各顶点的割前邻域中，则三个实际角不小于相应模型角；等价地，其中的最短铰链第三边不大于同腿同角模型的第三边。
\end{proposition}
\begin{proof}
 设三边长为 \(a,b,c\)。三角不等式给出
 \begin{equation}\label{g3:eq:triangle-radius-perimeter-bound}
 d(p,\sigma(t))\le\min\{c+t,b+a-t\}\le\frac{a+b+c}{2}<b_k.
 \end{equation}
 模型中同样成立。令 \(h(t)=m_k(d(p,\sigma(t)))\)，\(\bar h(t)=m_k(d_k(\bar p,\bar\sigma(t)))\)。式\eqref{g3:eq:modified-distance-hessian-smooth}给出 \(h''+kh\le1\)，模型给出 \(\bar h''+k\bar h=1\)。若边通过 \(p\)，\(m_k(d(p,\cdot))\) 在法坐标中是 \(m_k(|x|)\)，其偶次幂展开说明它在中心仍光滑，且 Hessian 为 \(g_p\)。故这不产生额外奇点。两函数的端点值相同，\(a<b_k\)，由\cref{g3:lem:one-dimensional-support-gluing}得到 \(h\ge\bar h\)。\(m_k\) 严格递增，给出距离比较。

 在顶点 \(q\) 处，两侧的距离相同，比较不等式的右导数给出
 \[
 -\cos\angle pqr
 =\left.\frac{\mathrm d}{\mathrm dt}d(p,\sigma(t))\right|_{0}
 \ge-\cos\widetilde\angle_k(pqr).
 \]
 这里使用距离第一变分公式。余弦在 \([0,\pi]\) 上递减，故实际角较大；交换三个顶点得到全部角比较。固定两腿后，式\eqref{g3:eq:model-hinge-cosine}把角比较转换为铰链第三边比较。
\end{proof}

\subsection{全局三角比较的证明}
\label{g3:subsec:0903:03}

全局化的关键是在距离不光滑的点仍构造式\eqref{g3:eq:one-dimensional-upper-support}中的支撑。取到该点的一条最短线，把起点向前移动一小段，再补上这段长度；三角不等式保证新函数在原距离上方。第七章的\cref{g3:lem:shifted-distance-smooth}已经证明这个新距离在接触点光滑。

\begin{proposition}\label{g3:prop:global-modified-distance-hessian}
 设 \((M^n,g)\) 连通、完备、光滑且无边界，\(n\ge2\)、\(k\in\mathbb R\)、\(K\ge k\)，并固定 \(p\in M\)。若单位速测地线 \(\sigma\) 的像位于 \(B(p,b_k)\)，则
 \[
 h(t)=m_k(d(p,\sigma(t)))
 \]
 在每个内点以上支撑意义满足 \(h''+kh\le1\)。
\end{proposition}
\begin{proof}
 固定参数 \(t_0\)，记 \(x=\sigma(t_0)\)、\(\rho=d(p,x)\)。若 \(x=p\)，前面法坐标中的偶函数计算给出 \(h''(t_0)=1\)，可以直接以 \(h\) 为局部支撑。

 若 \(\rho>0\)，由完备性取从 \(p\) 到 \(x\) 的单位速最短线 \(\eta:[0,\rho]\to M\)。对 \(0<\delta<\rho\)，令
 \[
 p_\delta=\eta(\delta),\qquad r_\delta(y)=\delta+d(p_\delta,y).
 \]
 则 \(r_\delta\ge d(p,\cdot)\)，且 \(r_\delta(x)=\rho\)。移基点引理保证它在 \(x\) 附近光滑。缩小邻域使 \(r_\delta<b_k\)，于是 \(H_\delta=m_k\circ r_\delta\circ\sigma\) 是 \(h\) 在 \(t_0\) 的光滑上支撑。

 记 \(z=\langle\dot\eta(\rho),\dot\sigma(t_0)\rangle\)，则 \(|z|\le1\)。从 \(p_\delta\) 到 \(x\) 的段长是 \(\rho-\delta<b_k\)，在割点以前；Hessian 比较给出
 \[
 \begin{split}
 H_\delta''(t_0)
 &\le c_k(\rho)z^2+s_k(\rho)
       \frac{c_k(\rho-\delta)}{s_k(\rho-\delta)}(1-z^2)\\
 &=c_k(\rho)+
   \frac{s_k(\delta)}{s_k(\rho-\delta)}(1-z^2).
 \end{split}
 \]
 最后一步使用 \(s_k(\rho)c_k(\rho-\delta)-c_k(\rho)s_k(\rho-\delta)=s_k(\delta)\)，可由正弦、双曲正弦的差角公式核验，\(k=0\) 时就是 \(\rho-(\rho-\delta)=\delta\)。因 \(0<\rho<b_k\)，末项随 \(\delta\downarrow0\) 趋于零。给定任意 \(\varepsilon>0\)，取足够小的 \(\delta\)，便有
 \[
 H_\delta''(t_0)+kh(t_0)\le1+\varepsilon.
 \]
 这正是所需的上支撑条件。
\end{proof}

这个构造只使用接触点处的一条最短线。不同参数处可以选择不同的最短线；即使这些选择不连续，上支撑拼接引理仍然适用，因为它的反证只在一个负最小点使用一个局部支撑。

\begin{theorem}[Toponogov 三角比较]\label{g3:thm:toponogov-triangle}
 设 \((M^n,g)\) 连通、完备、光滑且无边界，\(n\ge2\)、\(k\in\mathbb R\)，并且每个截面曲率满足 \(K\ge k\)。给定由三条选定最短测地线构成的三角形 \(pqr\)，设其三边长为 \(a=d(q,r)\)、\(b=d(p,r)\)、\(c=d(p,q)\)，并满足
 \[
 P=a+b+c<2b_k.
 \]
 令 \(\bar p\bar q\bar r\subset M_k^2\) 为比较三角形，\(\sigma:[0,a]\to M\) 及 \(\bar\sigma:[0,a]\to M_k^2\) 为对应的单位速边 \(qr\) 与 \(\bar q\bar r\)。则
 \begin{equation}\label{g3:eq:toponogov-vertex-side}
 d(p,\sigma(t))\ge d_k(\bar p,\bar\sigma(t))\qquad(0\le t\le a).
 \end{equation}
 三个实际角也分别不小于相应模型角。结论对所选的任意最短边成立，不要求这些最短边唯一，也不要求三角形位于某个法邻域中。
\end{theorem}
\begin{proof}
 由式\eqref{g3:eq:triangle-radius-perimeter-bound}，整条 \(\sigma\) 位于 \(B(p,b_k)\)，且 \(a<b_k\)。因此\cref{g3:prop:global-modified-distance-hessian}给出
 \[
 h=m_k\circ d(p,\cdot)\circ\sigma,
 \qquad h''+kh\le1
 \]
 的逐点上支撑不等式。模型函数 \(\bar h=m_k\circ d_k(\bar p,\cdot)\circ\bar\sigma\) 满足 \(\bar h''+k\bar h=1\)，两端分别取 \(m_k(c),m_k(b)\)。\cref{g3:lem:one-dimensional-support-gluing}遂给出 \(h\ge\bar h\)。在 \([0,b_k)\) 上用 \(m_k\) 的严格单调性，就得到式\eqref{g3:eq:toponogov-vertex-side}。这一步把局部 Hessian 控制拼成全局距离控制，已经包括 \(\sigma\) 经过割迹的情形。

 为比较 \(q\) 处的角，沿三角形中选定的边 \(pq\)，将起点 \(p\) 向前移动 \(\delta\)，构造上一命题中的 \(r_\delta\)。在 \(q\) 附近有
 \[
 r_\delta(\sigma(t))\ge d(p,\sigma(t))
 \ge d_k(\bar p,\bar\sigma(t)),
 \]
 且 \(t=0\) 时取等号。左端光滑，故其右导数不小于模型右导数。移基点没有改变所选边在 \(q\) 的终向量，距离第一变分给出
 \[
 -\cos\angle pqr\ge-\cos\widetilde\angle_k(pqr).
 \]
 因而 \(\angle pqr\ge\widetilde\angle_k(pqr)\)。模型边若通过 \(\bar p\)，模型距离在其内部会有尖点，但在端点 \(\bar q\ne\bar p\) 仍可求上述导数。循环交换顶点，得到三个角的结论。
\end{proof}

\begin{corollary}[最短铰链比较]\label{g3:cor:toponogov-hinge}
 设 \((M^n,g)\) 连通、完备、光滑且无边界，\(n\ge2\)、\(K\ge k\)。从 \(p\) 出发取两条最短段到 \(q,r\)，腿长为 \(c,b>0\)，夹角为 \(\alpha\)。若
 \(b+c+d(q,r)<2b_k\)，则同腿长、同夹角的模型铰链端点 \(\hat q,\hat r\) 满足
 \begin{equation}\label{g3:eq:toponogov-hinge}
 d(q,r)\le d_k(\hat q,\hat r).
 \end{equation}
 特别地，当 \(k=0\) 时，
 \[
 d(q,r)^2\le b^2+c^2-2bc\cos\alpha.
 \]
\end{corollary}
\begin{proof}
 若 \(q=r\)，结论直接成立。否则由完备性选取 \(qr\) 的一条最短段，应用\cref{g3:thm:toponogov-triangle}得到 \(\alpha\ge\widetilde\angle_k(qpr)\)。固定腿长时，模型第三边随夹角增加，故得到式\eqref{g3:eq:toponogov-hinge}。在正曲率情形，两腿都小于 \(b_k\)，比较铰链存在；即使其端点在特殊情形成为对径点，距离单调性仍由连续性成立。
\end{proof}

\subsection{正曲率模型中的周长限制}
\label{g3:subsec:0903:04}

当 \(k\le0\) 时，\(b_k=\infty\)，上述全局结果没有周长限制。当 \(k>0\) 时，条件 \(P<2b_k\) 同时完成三个任务：它使球面比较三角形由边长确定；使每条边长小于 \(b_k\)，从而一维边值比较中的正权函数存在；还使顶点到整条对边的距离严格小于 \(b_k\)，从而 \(m_k\) 严格递增，移基点误差中的分母 \(s_k(\rho-\delta)\) 不会在极限中消失。

\ProseParagraphBreak

这些限制可以在单位球面上直接看见。若三点把一条大圆分成三条长度都小于 \(\pi\) 的弧，则三条最短边的周长为 \(2\pi\)。在一条对边上，可能出现另一顶点的对径点；此时该距离等于 \(\pi\)，\(s_1(\pi)=0\)。若一条边的长度已经等于 \(\pi\)，连接其端点的最短半大圆也不唯一。边长与角的光滑对应、支撑估计和边值方程都不能在这些位置照搬严格范围内的除法。

\begin{example}\label{g3:exm:spherical-perimeter-threshold}
 在曲率 \(k>0\) 的模型球面上，取三条边都为 \(a\) 的非退化三角形。令 \(x=\sqrt{k}a\)，余弦公式给出
 \[
 \cos\alpha=\frac{\cos x}{1+\cos x}.
 \]
 右侧属于 \((-1,1)\) 的范围恰为 \(0<x<2\pi/3\)，即 \(3a<2b_k\)。当 \(x\uparrow2\pi/3\) 时，\(\alpha\uparrow\pi\)，三角形趋向把一条大圆三等分的退化构形。若取 \(x=3\pi/4\)，三边都小于 \(\pi\)，并满足普通三角不等式，但右侧小于 \(-1\)，因而不存在这样的球面比较三角形。这说明仅列每边小于 \(b_k\)，还不足以保证给定三边有模型。
\end{example}

在 \(K\ge k>0\) 的完备流形上，\cref{g3:thm:bonnet-myers}另外给出直径不超过 \(b_k\)。直径界是关于每一对点的结论，不能在选取比较三角形时替代所写的严格周长条件。本节所证明的版本始终在 \(P<2b_k\) 内使用；涉及临界周长时，应重新处理退化模型及对径点。

\subsection{等号和刚性情形}
\label{g3:subsec:0903:05}

在模型空间中当然处处取等号。反过来，等号首先控制三角形所扫过的方向。下面先证明一维等号的传播，再在径向扇形没有割点的条件下恢复它的二维度量。

\begin{proposition}\label{g3:prop:toponogov-profile-rigidity}
 设 \((M^n,g)\) 连通、完备、光滑且无边界，\(n\ge2\)、\(K\ge k\)，\(pqr\) 为周长小于 \(2b_k\) 的最短测地三角形。沿对应的单位速边 \(\sigma:[0,a]\to M\)、\(\bar\sigma:[0,a]\to M_k^2\)，若存在 \(t_0\in(0,a)\) 使
 \[
 d(p,\sigma(t_0))=d_k(\bar p,\bar\sigma(t_0)),
 \]
 则同一等式对所有 \(t\in[0,a]\) 成立。
\end{proposition}
\begin{proof}
 令 \(w=m_k(d(p,\sigma))-m_k(d_k(\bar p,\bar\sigma))\)。全局比较说明 \(w\ge0\)，而支撑构造说明 \(w''+kw\le0\) 以上支撑意义成立。固定任意 \(s<t_0<t\)，在 \([s,t]\) 上定义
 \[
 y(x)=\frac{s_k(t-x)}{s_k(t-s)}w(s)
      +\frac{s_k(x-s)}{s_k(t-s)}w(t).
 \]
 则 \(y''+ky=0\)，且与 \(w\) 有相同端点。因为 \(t-s<a<b_k\)，分母及内点的两个系数均为正。一维支撑拼接引理的齐次形式给出 \(w\ge y\)：其证明只需将常数一换成零，或直接对 \(w-y\) 使用该引理证明中的正权最小值论证。在 \(t_0\) 处有
 \(0=w(t_0)\ge y(t_0)\ge0\)，故 \(w(s)=w(t)=0\)。任意选择 \(s,t\) 即得 \(w\equiv0\)，再由 \(m_k\) 的单调性回到距离。
\end{proof}

\begin{proposition}[无割点扇形中的等号刚性]\label{g3:prop:toponogov-fan-rigidity}
 设 \((M^n,g)\) 连通、完备、光滑且无边界，\(n\ge2\)、\(K\ge k\)。设 \(pqr\) 是周长小于 \(2b_k\) 的非退化最短三角形，\(\sigma:[0,a]\to M\) 为边 \(qr\)，且
 \[
 \sigma([0,a])\cap\bigl(\operatorname{Cut}(p)\cup\{p\}\bigr)=\varnothing.
 \]
 若顶点到对边的比较在某个内点取等号，则从 \(p\) 到 \(\sigma(t)\) 的唯一最短段扫出的闭三角域，与 \(M_k^2\) 中的闭比较三角域具有相同的诱导黎曼度量。对应保持顶点及边上的弧长参数；三角域内部是一张全测地曲面。
\end{proposition}
\begin{proof}
 由\cref{g3:prop:toponogov-profile-rigidity}，记
 \[
 \rho(t)=d(p,\sigma(t))=d_k(\bar p,\bar\sigma(t)),\qquad
 u(t)=\frac{\log_p\sigma(t)}{\rho(t)},
 \]
 则 \(\rho,u\) 都光滑。这里 \(\log_p\) 是割前域上的指数映射逆。考虑径向参数化
 \[
 F(r,t)=\exp_p(ru(t)),\qquad 0<r\le\rho(t).
 \]
 我们先证明其诱导度量与模型相同，再排除扇形自交并证明全测地性。

 令 \(T\) 为径向单位向量，并在外边 \(\sigma(t)\) 上分解
 \[
 \dot\sigma(t)=\rho'(t)T+v(t),\qquad v(t)\perp T,
 \qquad |v(t)|^2=1-\rho'(t)^2.
 \]
 非退化模型中，边 \(\bar q\bar r\) 不可能在某点与来自 \(\bar p\) 的径向线相切；否则测地线初值唯一性会使这条边落在同一条过 \(\bar p\) 的测地线上。因此 \(|v(t)|>0\)。由于整个距离轮廓等于模型，沿 \(\sigma\) 对 \(m_k(\rho)\) 求二阶导数得到
 \[
 c_k(\rho)
 =c_k(\rho)\rho'^2+s_k(\rho)
       \operatorname{Hess}r_p(v,v).
 \]
 从而
 \[
 \operatorname{Hess}r_p(v,v)=\operatorname{ct}_k(\rho)|v|^2.
 \]
 距离仍在割前，\cref{g3:thm:hessian-comparison}的方向等号结论可用。沿固定 \(t\) 的径向线，角向 Jacobi 场 \(J=\partial_tF\) 满足
 \begin{equation}\label{g3:eq:triangle-equality-jacobi}
 J(r)=\frac{s_k(r)}{s_k(\rho(t))}P_{\rho(t)\to r}v(t).
 \end{equation}
 这里 \(J(0)=0\)，而由 \(\sigma(t)=F(\rho(t),t)\) 有 \(J(\rho(t))=v(t)\)，所以正是该方向等号中的边值场。

 在模型中选极角 \(\vartheta\)，使 \(\bar\sigma\) 的极角随 \(t\) 增加。它从零增至模型顶角 \(\bar\alpha\in(0,\pi)\)，并满足
 \[
 \vartheta'(t)=\frac{\sqrt{1-\rho'(t)^2}}{s_k(\rho(t))}>0.
 \]
 可在球面或双曲面线性模型中核对这一选择：模型边点是两个端点的正线性组合，系数为
 \(s_k(a-t)/s_k(a)\) 与 \(s_k(t)/s_k(a)\)，其径向方向扫过夹角小于 \(\pi\) 的扇形；平面中则是普通线段插值。以 \(\vartheta\) 代替 \(t\)，Gauss 引理及式\eqref{g3:eq:triangle-equality-jacobi}给出
 \begin{equation}\label{g3:eq:triangle-equality-polar-metric}
 F^*g=\mathrm dr^2+s_k(r)^2\,\mathrm d\vartheta^2.
 \end{equation}
 因为角向场非零，这在 \(r>0\) 处已经是一张浸入曲面的度量。

 以下说明该曲面没有额外的弯曲。记
 \(X=\partial_r\)、\(Y=s_k(r)^{-1}\partial_\vartheta\)，它们是诱导度量的正交单位标架。径向线是环境测地线，而式\eqref{g3:eq:triangle-equality-jacobi}说明 \(Y\) 沿径向平行，故曲面的第二基本形式满足
 \[
 B(X,X)=B(X,Y)=0.
 \]
 Hessian 等号同时给出环境径向截面曲率 \(K(X,Y)=k\)。若 \(N\) 是单位法向量，考虑二维平面
 \(\operatorname{span}\{\cos s\,X+\sin s\,N,Y\}\)。它的截面曲率在 \(s=0\) 达到下界 \(k\)，对 \(s\) 求导得到
 \[
 0=2\langle R(X,Y)Y,N\rangle.
 \]
 因此 \((R(X,Y)Y)^\perp=0\)。余维为零时这一结论直接成立。

 由式\eqref{g3:eq:triangle-equality-polar-metric}，曲面自身的联络满足
 \[
 \nabla_XY=0,\qquad \nabla_YX=\operatorname{ct}_k(r)Y,
 \qquad \nabla_YY=-\operatorname{ct}_k(r)X.
 \]
 将这些关系代入\cref{g3:thm:codazzi-equation}，令 \(C=B(Y,Y)\)，得到
 \[
 0=(\nabla B)(X;Y,Y)-(\nabla B)(Y;X,Y)
   =\nabla_X^\perp C+\operatorname{ct}_k(r)C.
 \]
 即 \(\nabla_X^\perp(s_k(r)C)=0\)。在外边上，\(\sigma\) 是环境测地线，故 \(B(\dot\sigma,\dot\sigma)=0\)。其角向分量非零，而另外两项 \(B(X,X),B(X,Y)\) 已为零，所以 \(C=0\) 于外边。沿每条径向线解上述线性方程，得 \(C=0\) 于全扇形。因此其内部全测地。

 最后检验参数化的单射性和中心的光滑性。将 \(u\) 也用 \(\vartheta\) 参数化，式\eqref{g3:eq:triangle-equality-jacobi}在 \(r\downarrow0\) 的极限给出 \(|u'|=1\)。由于第二基本形式为零，环境中有
 \[
 \overline\nabla_{\partial_\vartheta}\partial_\vartheta F
   =-s_k(r)c_k(r)\,\partial_rF.
 \]
 在 \(p\) 的法坐标中 \(F(r,\vartheta)=ru(\vartheta)\)，且 \(\Gamma(x)=O(|x|)\)。故上式的坐标形式为
 \[
 r u''(\vartheta)+O(r^3)
   =-s_k(r)c_k(r)u(\vartheta)=-r u(\vartheta)+O(r^3).
 \]
 除以 \(r\) 并令 \(r\downarrow0\)，得到 \(u''=-u\)。于是 \(u\) 是 \(S_pM\) 上的一段单位速大圆，且其参数长度 \(\bar\alpha<\pi\)，所以没有重复方向。整个径向扇形位于 \(p\) 的割前域，指数映射在那里单射。因此 \(F\) 在三角域上没有自交。

 取把模型起点切平面映到 \(\operatorname{span}\{u(0),u'(0)\}\) 的线性等距，模型三角域到所得扇形的映射就是这个线性等距与两边指数坐标的复合。它在中心光滑，微分保持度量；在其余点保持度量由式\eqref{g3:eq:triangle-equality-polar-metric}给出。它也把三条模型边映到给定三角形的三条边，完成所述刚性。
\end{proof}

这里得到的是三角域的诱导度量及其全测地性。环境中未被扇形扫过的二维平面仍可具有其他曲率。例如 \(\mathbb R^2\times S^2\) 的平面因子内，每个三角形都达到 \(k=0\) 的比较等号，而球面因子的截面曲率为正。若只知某一点对的距离等号而没有上述条件，也不能把它扩大解释为一个环境开集的常曲率性。
% 纲要条目排为编号小节标题，并进入目录与书签；说明段仍保留于纲要注释中。

% !TeX root = ../../Book3.tex
% 纲要标题：### 9.4 三角比较的几何后果
% 纲要位置：几何学计划纲要.md，第 2452 行。
% 正文状态：初稿；尚未正式审核。

% 纲要内容开始
% 
% * 凸集、测地三角形和投影
% * 球面型直径约束
% * CAT 上界与 Alexandrov 下界的模型对照
% 
% ---
% 
% 纲要内容结束

\section{三角比较的几何后果}
\label{g3:sec:0904}

三角比较把曲率条件转成有限个点之间的距离关系。非正曲率配合单连通性，给出稳定的凸性与最近点投影；正曲率下界则使最短线长度受到球面模型的限制。我们先分别推导这些结论，最后用度量三角形的语言比较上下曲率界。前两小节及 CAT(0) 的证明只使用本章常规部分；最后关于下三角比较的结论使用\secref{g3:sec:0903}。

\subsection{凸集、测地三角形和投影}
\label{g3:subsec:0904:01}

沿用\cref{g3:def:hadamard-manifold}的 Hadamard 流形：它连通、完备、单连通、光滑且无边界，截面曲率非正。任意两点之间的测地线段唯一，因而“包含两点之间的线段”不再涉及最短线的选择。

前面已经用过的\term[cedi tuxing]{测地凸}[geodesically convex]集合 \(C\subset M\)，就是包含其任意两点之间唯一最短线段的集合。对于这样的集合，垂足可以用最小化距离定义，即使集合的边界并不光滑。

距离平方的强凸性有一个方便的有限距离表达。设 \(p,x,y\in M\)，\(\gamma:[0,1]\to M\) 是从 \(x\) 到 \(y\) 的常速最短线，\(a=d(x,y)\)。由\cref{g3:prop:hadamard-squared-distance-convex}，
\[
	\frac{\mathrm d^2}{\mathrm dt^2}\frac12d(p,\gamma(t))^2
	\ge a^2.
\]
减去 \(a^2t^2/2\) 后得到凸函数，再用端点连线在其图像上方，便得
\begin{equation}\label{g3:eq:hadamard-cn}
	d(p,\gamma(t))^2
	\le(1-t)d(p,x)^2+td(p,y)^2-t(1-t)d(x,y)^2.
\end{equation}
在 Euclid 空间中，展开平方便知这里恒为等式。非正曲率保留的方向是“线段内点的距离平方不超过这个模型值”。

\begin{theorem}[闭凸集的最近点投影]\label{g3:thm:hadamard-projection}
	设 \(M\) 为有限维 Hadamard 流形，\(C\subset M\) 非空、闭且测地凸。对每个 \(x\in M\)，存在唯一 \(u=P_C(x)\in C\) 使
	\[
		d(x,u)=\inf_{q\in C}d(x,q).
	\]
	该点由下列条件刻画：
	\begin{equation}\label{g3:eq:projection-angle}
		\langle\log_u x,\log_u q\rangle\le0
		\qquad(q\in C).
	\end{equation}
	此外，对所有 \(q\in C\)，
	\begin{equation}\label{g3:eq:projection-pythagoras}
		d(x,q)^2\ge d(x,P_Cx)^2+d(P_Cx,q)^2.
	\end{equation}
	映射 \(P_C:M\to C\) 为1-Lipschitz。
\end{theorem}
\begin{proof}
	取 \(q_0\in C\)。一列逼近距离下确界的点最终落在闭球 \(\overline B(x,d(x,q_0)+1)\) 内；该球由完备性和 Hopf--Rinow 定理为紧集。取收敛子列，闭性保证极限仍在 \(C\)，距离连续性说明极限达到下确界。若不同的 \(u,v\) 都达到最小值 \(\rho\)，则式\eqref{g3:eq:hadamard-cn}在其中点 \(m\in C\) 处给出
	\[
		d(x,m)^2\le\rho^2-\tfrac14d(u,v)^2<\rho^2,
	\]
	矛盾，故最近点唯一。

	令 \(\eta(t)=\exp_u(t\log_u q)\)。当 \(u=P_Cx\) 时，函数 \(f(t)=\tfrac12d(x,\eta(t))^2\) 在 \(t=0\) 的右导数非负。梯度公式给出
	\[
		f'(0)=-\langle\log_u x,\log_u q\rangle\ge0.
	\]
	反过来，若 \(u\in C\) 满足式\eqref{g3:eq:projection-angle}，则 \(f''(t)\ge d(u,q)^2\)，积分两次得到
	\[
		\tfrac12d(x,q)^2
		\ge\tfrac12d(x,u)^2+f'(0)+\tfrac12d(u,q)^2
		\ge\tfrac12d(x,u)^2+\tfrac12d(u,q)^2.
	\]
	因此 \(u\) 就是最近点，并同时得到式\eqref{g3:eq:projection-pythagoras}。对不等于 \(u\) 的 \(x,q\)，条件\eqref{g3:eq:projection-angle}意味着从 \(u\) 指向 \(x\) 与指向 \(q\) 的夹角至少为直角。

	最后证明投影不增加距离。设 \(u=P_Cx\)、\(v=P_Cy\)，先考虑 \(u\ne v\)。取
	\[
		\alpha(s)=\exp_u(s\log_u x),\qquad
		\beta(s)=\exp_v(s\log_v y),\qquad 0\le s\le1,
	\]
	并令 \(E(s)=\tfrac12d(\alpha(s),\beta(s))^2\)。\Cref{g3:prop:hadamard-geodesic-uniqueness}已经证明唯一测地线及 \(\log\) 映射光滑依赖两个端点。因此 \(E\) 在两端点相合处仍光滑。

	用 \([0,1]\) 上的常速最短线连接 \(\alpha(s)\) 与 \(\beta(s)\)，得到光滑变分。两条端点轨迹都是测地线，故\cref{g3:thm:energy-second-variation}的端点加速度项为零；又因 \(K\le0\)，其积分项非负。因此 \(E''(s)\ge0\)。这个推导在端点相合时也成立，此时中间测地线为常曲线。若 \(T\) 是 \(s=0\) 时从 \(u\) 到 \(v\) 的常速切向场，第一变分给出
	\[
		E'(0)=\langle\log_v y,T(1)\rangle
		-\langle\log_u x,T(0)\rangle\ge0.
	\]
	这里 \(T(0)=\log_u v\)、\(T(1)=-\log_v u\)，两项的符号分别来自 \(v,u\) 的投影角条件。于是 \(E(1)\ge E(0)\)，即 \(d(x,y)\ge d(u,v)\)。若 \(u=v\)，此不等式直接成立。
\end{proof}

非空闭测地凸集因而具有连续的收缩映射 \(P_C\)，在 \(C\) 上等于恒等。由最短线定义
\[
	H(x,t)=\exp_x\bigl(t\log_x(P_Cx)\bigr)
\]
还得到从恒等映射到 \(P_C\) 的连续形变，且在 \(C\) 上每点不动。闭性负责取得最小值，凸性负责唯一性及投影角；若删除凸性，Euclid 平面中两个分离点的中垂线上便有两个最近点。

\subsection{球面型直径约束}
\label{g3:subsec:0904:02}

非正曲率在单连通覆盖中容许测地线无限保持最短，统一的正截面曲率下界产生相反效果。这里可以直接接上第八章的直径估计。

\begin{corollary}\label{g3:cor:positive-sectional-diameter}
	设 \((M^n,g)\) 连通、完备、光滑且无边界，\(n\ge2\)。若所有截面曲率满足 \(K\ge k>0\)，则
	\[
		\operatorname{diam}M\le b_k=\frac{\pi}{\sqrt k},
	\]
	且 \(M\) 紧、基本群有限。
\end{corollary}
\begin{proof}
	在任一单位向量 \(T\) 的垂直补中取正交单位基 \(e_1,\ldots,e_{n-1}\)，得到
	\[
		\operatorname{Ric}(T,T)=\sum_{i=1}^{n-1}K(T\wedge e_i)\ge(n-1)k.
	\]
	\Cref{g3:thm:bonnet-myers,g3:cor:myers-finite-fundamental-group}的条件全部满足，故得到所述结论。
\end{proof}

半径 \(1/\sqrt k\) 的圆球面达到直径界。相同的局部曲率也可以产生较小的整体直径：在实射影空间的标准商度量中，若 \(x,y\) 为半径 \(a=1/\sqrt k\) 球面上的代表元，则
\[
	d([x],[y])=\min\{d_{S^n(a)}(x,y),d_{S^n(a)}(x,-y)\}.
\]
这来自覆盖道路的两个可能终点。球面上这两个距离之和为 \(\pi a\)，所以射影空间直径为 \(\pi a/2\)，且当 \(x\perp y\) 时达到。局部曲率不区分这两个空间，整体的最短线则会看到商的识别。

\ProseParagraphBreak

对正曲率模型，直径范围还影响三角比较的陈述。一个球面模型的短边三角形要求周长严格小于 \(2b_k\)；这个条件保证每边短于 \(b_k\)，从而避开对径点处的最短线不唯一。直径上界本身只逐边给出 \(a,b,c\le b_k\)，不能把三边各自的上界相加后就声称已经满足三角比较的严格周长条件。\secref{g3:sec:0903}所证明的版本将该条件明确列在定理中。

边界情形在标准球面上就能看到：取一对对径点，二者之间有无穷多条长度 \(b_k\) 的最短半大圆；再取其中一条的内点，三个距离的总和恰为 \(2b_k\)。因此在这样的端点配置中，单靠三边长无法选定唯一的模型短边三角形。另一方面，曲率上界 \(K\le k\) 不产生上述有限直径结论：Euclid 空间对任何 \(k>0\) 都满足这个上界，却有无限直径。

\subsection{CAT 上界与 Alexandrov 下界的模型对照}
\label{g3:subsec:0904:03}

若空间中任意两点可以由长度等于其距离的路径连接，并把这样的路径按常速参数化，就称它为测地度量空间。黎曼完备性保证连通流形具有这一性质。三角比较的条件只涉及这些路径上的点与距离，因此可以脱离曲率张量表达。

\begin{definition}\label{g3:def:cat-comparison}
	设 \(X\) 是测地度量空间，\(k\in\mathbb R\)。对三边均为选定最短线、周长小于 \(2b_k\) 的三角形 \(pqr\)，在常曲率 \(k\) 的标准二维模型中取同边长三角形 \(\bar p\bar q\bar r\)。边上的对应点以距相应端点的比例确定。若任意这样的三角形及其边上任意两点 \(x,y\) 均满足
	\begin{equation}\label{g3:eq:cat-upper-distance}
		d_X(x,y)\le d_k(\bar x,\bar y),
	\end{equation}
	则称 \(X\) 为 \(\mathrm{CAT}(k)\) 空间。允许退化三角形，并以连续极限理解对应点。这里 \(b_k=\infty\) 当 \(k\le0\)。
\end{definition}
\begin{figure}[htbp]
	\centering
	% !TeX root = ../Book3.tex
\begin{tikzpicture}[x=.85cm,y=.85cm,>=stealth]
	\coordinate (p) at (0,2.3);
	\coordinate (q) at (-1.65,0);
	\coordinate (r) at (1.65,0);
	\coordinate (z) at (.35,0);
	\draw[thick] (q) to[bend left=9] node[left] {\(c\)} (p);
	\draw[thick] (p) to[bend left=9] node[right] {\(b\)} (r);
	\draw[thick] (q)--(r);
	\draw[dashed] (p)--(z);
	\foreach \a in {p,q,r,z} {\fill (\a) circle (1.5pt);}
	\node[above] at (p) {\(p\)};
	\node[below left] at (q) {\(q\)};
	\node[below right] at (r) {\(r\)};
	\node[below] at (z) {\(z\)};
	\node[below] at (-.7,-.3) {\(ta\)};
	\node[below] at (1,-.3) {\((1-t)a\)};
	\draw[->] (2.2,1.15)--(3.3,1.15) node[midway,above] {对应};
	\begin{scope}[xshift=5.5cm]
		\coordinate (P) at (0,2.3);
		\coordinate (Q) at (-1.65,0);
		\coordinate (R) at (1.65,0);
		\coordinate (Z) at (.35,0);
		\draw[thick] (Q)--node[left] {\(c\)} (P)--node[right] {\(b\)} (R)--(Q);
		\draw[dashed] (P)--(Z);
		\foreach \a in {P,Q,R,Z} {\fill (\a) circle (1.5pt);}
		\node[above] at (P) {\(\bar p\)};
		\node[below left] at (Q) {\(\bar q\)};
		\node[below right] at (R) {\(\bar r\)};
		\node[below] at (Z) {\(\bar z\)};
		\node[below] at (-.7,-.3) {\(ta\)};
		\node[below] at (1,-.3) {\((1-t)a\)};
	\end{scope}
\end{tikzpicture}

	\caption{三角形与同边长模型的对应关系。底边长为 \(a\)，对应点 \(z,\bar z\) 距左端点都为 \(ta\)；虚线所代表的距离是顶点到对边比较的对象。右图以 Euclid 模型示意。}
	\label{g3:fig:triangle-correspondence}
\end{figure}
如\cref{g3:fig:triangle-correspondence}所示，比较时先固定三条边长，再按边上的距离比例确定对应点。特别地，顶点到对边的比较也包含在式\eqref{g3:eq:cat-upper-distance}中。
\term[CAT kongjian]{CAT 空间}[CAT space]的名称汇集了 Cartan、Alexandrov 与 Toponogov 的姓氏。Alexandrov 的工作使曲率比较进入了没有光滑坐标的度量空间。这里的 Alexandrov 是 Aleksandr Danilovich Alexandrov。\footnote{Aleksandr Danilovich Alexandrov（1912--1999），苏联、俄罗斯数学家，研究凸曲面、整体几何和度量空间中的曲率比较；与拓扑学家 Pavel Sergeevich Alexandrov 为不同人物。}

\begin{theorem}\label{g3:thm:hadamard-cat-zero}
	有限维 Hadamard 流形连同其黎曼距离是 \(\mathrm{CAT}(0)\) 空间。
\end{theorem}
\begin{proof}
	由完备性它是测地度量空间。取三角形 \(pqr\)，并令 \(x\in[p,q]\)、\(y\in[p,r]\)，写
	\[
		d(p,x)=s\,d(p,q),\qquad d(p,y)=t\,d(p,r),
		\qquad 0\le s,t\le1.
	\]
	在式\eqref{g3:eq:hadamard-cn}中先以 \(r\) 为基点、\([p,q]\) 为线段，得到
	\[
		d(r,x)^2
		\le(1-s)d(r,p)^2+s\,d(r,q)^2-s(1-s)d(p,q)^2
		=|\bar r-\bar x|^2.
	\]
	最后的等号是 Euclid 空间的平方展开。再以 \(x\) 为基点、\([p,r]\) 为线段，得到
	\[
	\begin{aligned}
		d(x,y)^2
		&\le(1-t)d(x,p)^2+t\,d(x,r)^2-t(1-t)d(p,r)^2\\
		&\le(1-t)|\bar x-\bar p|^2+t|\bar x-\bar r|^2
		-t(1-t)|\bar p-\bar r|^2\\
		&=|\bar x-\bar y|^2.
	\end{aligned}
	\]
	同一条边上的对应点距离直接相等；不同边总有一个公共顶点，重新标记即可应用上述证明。因此所有边点对都满足式\eqref{g3:eq:cat-upper-distance}。
\end{proof}

这个结果同时解释一个看似相反的比较方向。\cref{g3:prop:hadamard-exponential-expansion}在固定初始夹角与两条径向边长时，使第三边不小于 Euclid 铰链；CAT(0) 则固定全部三边长，比较边内点之间的距离。两者固定的数据不同：在同边长三角形中，非正曲率使顶角不超过 Euclid 模型角，边间距因而满足上述上界。

\ProseParagraphBreak

下曲率界对应另一方向。对同边长模型三角形，要求每个顶点 \(p\) 与对边上的对应点 \(z,\bar z\) 满足
\begin{equation}\label{g3:eq:alexandrov-lower-distance}
	d(p,z)\ge d_k(\bar p,\bar z)
\end{equation}
是 Alexandrov 下曲率界的三角比较形式。\secref{g3:sec:0903}在连通完备光滑无边界流形的 \(K\ge k\) 条件下，完整证明了周长小于 \(2b_k\) 的这一结论。与 CAT 上比较不同，这个整体下比较不要求流形单连通。

\begin{table}[htbp]
	\centering
	\caption{以同边长模型为参照的两种比较。下行采用本章证明的顶点到对边形式。}
	\label{g3:tab:triangle-comparison-directions}
	\begin{tabularx}{\textwidth}{@{}lXX@{}}
		\toprule
		类型 & 本章已证的光滑情形 & 对应距离的方向\\
		\midrule
		CAT(0) 上比较 & 完备、单连通，\(K\le0\) & 边点间距不超过 Euclid 模型\\
		\(k\) 下三角比较 & 完备，\(K\ge k\)，周长 \(<2b_k\) & 顶点到对边距离不小于模型\\
		\bottomrule
	\end{tabularx}
\end{table}

局部曲率上界与整体 CAT 条件之间还隔着拓扑。例如平坦二维环面 \((\mathbb R/\ell\mathbb Z)\times(\mathbb R/\ell\mathbb Z)\) 的曲率恒为零。沿第一圆周取三个等间距点 \(p,q,r\)，以三段长 \(\ell/3\) 的短弧为边；令 \(z\) 为不经过 \(p\) 的边 \([q,r]\) 的中点。则
\[
	d(p,z)=\ell/2,\qquad
	|\bar p-\bar z|=\frac{\sqrt3}{2}\frac\ell3
	=\frac{\ell}{2\sqrt3}.
\]
它违反 CAT(0) 不等式，虽然每个足够小的坐标邻域都是 Euclid 的。其万有覆盖是 Euclid 平面，才恢复整体 CAT(0) 比较。曲率的局部条件、最短线的整体选择与覆盖的拓扑，在这个例子中分别可见。
% 纲要条目排为编号小节标题，并进入目录与书签；说明段仍保留于纲要注释中。


\begin{chaptersummary}
	非正曲率先排除 Jacobi 场的再次归零，并使指数映射的微分不缩短向量。拉回度量因此完备，覆盖论证把局部可逆性提升为整体结构；单连通性再使覆盖只有一层。由此得到全局唯一最短线、距离平方的强凸性，以及闭凸集的唯一非扩张投影。

	\ProseParagraphBreak

	常曲率的角向 Jacobi 场精确等于 \(s_k\) 乘平行场，决定了局部极坐标度量。正曲率时，公共终点和反向末速度共同保证模型南极处的光滑延拓。完整的单连通模型分类随后把一般空间形式归结为自由、适当不连续的等距群商；群在模型等距群中的共轭类决定商的等距类。

	\ProseParagraphBreak

	Toponogov 的距离比较由局部 Hessian 上界出发，以移动基点构造割点处的光滑上支撑，再经一维边值比较传播到整个区间。模型首零尺度规定了正曲率的周长范围，等号则需要逐项追踪支撑及径向曲率中的损失。CAT 上比较与下三角比较固定同样的边长而取相反的距离方向。下一章将以长度空间与 Gromov--Hausdorff 距离组织这些度量信息，研究空间变化时可以保留什么。
\end{chaptersummary}

\BookExercises

\begin{exercise}\label{g3:ex:hyperbolic-exp-unbounded}
	在曲率 \(-1\) 的双曲空间 \(\mathbb H^n\)、\(n\ge2\) 中固定 \(p\)。证明 \(\exp_p\) 虽满足全局距离扩张，却不是从 Euclid 切空间到 \(\mathbb H^n\) 的全局 Lipschitz 映射。求其在半径 \(R\) 的切球上的微分算子范数上界，并说明该界是最优的上确界。
\end{exercise}
\begin{solution}
	取单位向量 \(u\) 及 \(w\perp u\)。常曲率 Jacobi 场给出
	\[
		(\mathrm d\exp_p)_{ru}(au+w)
		=a\dot\gamma_u(r)+\frac{\sinh r}{r}P_rw.
	\]
	径向与横向分量正交，故微分范数为 \(\sinh r/r\)，在 \(r=0\) 处取极限1。其导数的分子为 \(r\cosh r-\sinh r\)，这个分子在零处为零、导数为 \(r\sinh r>0\)，因此范数随 \(r>0\) 严格增加。半径 \(R\) 的开切球上的上确界为 \(\sinh R/R\)，可由趋于边界的横向单位向量逼近。沿切球中的直线段积分微分界，得到该球上的同一 Lipschitz 常数。

	若全域映射存在 Lipschitz 常数 \(L\)，将距离不等式用于 \(ru\) 与 \(ru+tw\)，除以 \(|t|\) 并令 \(t\to0\)，便有 \(\sinh r/r\le L\) 对所有 \(r\) 成立。这与该比值趋于无穷矛盾。全局扩张控制下界；横向微分随半径增长说明上界可以完全不同。
\end{solution}

\begin{exercise}\label{g3:ex:flat-torus-exponential-cover}
	在标准平坦环面 \(T^n=\mathbb R^n/\mathbb Z^n\)、\(n\ge2\) 上，取 \(p=[0]\) 并把 \(\mathrm T_pT^n\) 识别为 \(\mathbb R^n\)。
	\begin{enumerate}
		\item 求 \(\exp_p\)，证明其微分处处可逆而它不是单射。
		\item 对 \(q=[x]\)，用格点给出所有定义于 \([0,1]\) 的从 \(p\) 到 \(q\) 的仿射参数测地线，并辨认哪些最短。
		\item 取 \(x=(1/2,0,\ldots,0)\)，计算最短线的条数，并解释这不与非正曲率的无共轭点结论矛盾。
	\end{enumerate}
\end{exercise}
\begin{solution}
	覆盖投影局部等距，Euclid 直线投影后满足测地线方程，因此
	\(\exp_p(v)=[v]\)。其微分在局部覆盖坐标中为恒等映射，故处处可逆；但 \(v\) 与 \(v+m\)、\(m\in\mathbb Z^n\) 有相同像。

	从 \(p\) 出发的测地线提升为从零出发的直线，故恰为
	\[
		\gamma_m(t)=[t(x+m)],\qquad m\in\mathbb Z^n.
	\]
	长度为 \(|x+m|\)。任意连接道路提升后的终点也为某个 \(x+m\)，其长度至少为这个 Euclid 距离；反向以线段达到。因此
	\(d(p,q)=\min_m|x+m|\)，并且上述线段最短当且仅当 \(m\) 达到这个最小值。格点离散且范数在无穷远趋于无穷，最小值确实达到。

	所给 \(x\) 的最小值为 \(1/2\)，只有 \(m=0\) 和 \(m=(-1,0,\ldots,0)\) 达到，故恰有两条最短线。指数映射的不同初向量在此被整体识别，但微分没有退化；割点由最短线竞争造成，并不一定为共轭点。这里缺少 Cartan--Hadamard 的单连通条件。
\end{solution}

\begin{exercise}\label{g3:ex:hadamard-busemann-concavity}
	设 \(M\) 为 Hadamard 流形，\(\gamma\) 为单位速测地射线。沿用第八章
	\[
		b_\gamma(x)=\lim_{T\to\infty}\bigl(T-d(x,\gamma(T))\bigr)
	\]
	的符号，证明 \(b_\gamma\) 沿测地线凹，且对每个实数 \(a\)，集合 \(\{b_\gamma\ge a\}\) 为非空闭测地凸集。进一步说明它有全域的1-Lipschitz 最近点投影。
\end{exercise}
\begin{solution}
	对固定 \(T\)，距离 \(d(\,\cdot\,,\gamma(T))\) 由\cref{g3:cor:hadamard-distance-balls-convex}为测地凸函数，所以 \(b_T=T-d(\,\cdot\,,\gamma(T))\) 为凹函数。若 \(\eta:[0,1]\to M\) 是常速测地线段，则
	\[
		b_T(\eta(t))\ge(1-t)b_T(\eta(0))+tb_T(\eta(1)).
	\]
	第八章已经证明逐点极限存在，令 \(T\to\infty\) 即保留该不等式。两个端点的值若都至少为 \(a\)，线段上所有值也至少为 \(a\)，故超水平集测地凸。Busemann 函数为1-Lipschitz，因而连续，超水平集闭；沿射线有 \(b_\gamma(\gamma(s))=s\)，取 \(s\ge\max\{a,0\}\) 即知其非空。最后应用\cref{g3:thm:hadamard-projection}得到唯一、1-Lipschitz 的最近点投影。符号若改为 \(d-T\)，相应结论应改写成凸函数的次水平集。
\end{solution}

\begin{exercise}\label{g3:ex:flat-tori-isometry-lattices}
	设 \(\Lambda,\Lambda'\subset\mathbb R^n\) 为满秩格。证明平坦环面 \(T_\Lambda\) 与 \(T_{\Lambda'}\) 等距，当且仅当存在 \(Q\in O(n)\) 使 \(Q\Lambda=\Lambda'\)。在二维构造面积相同而不等距的平坦环面。
\end{exercise}
\begin{solution}
	Euclid 空间的黎曼等距必为 \(A(x)=Qx+b\)：令 \(b=A(0)\)、\(Q=\mathrm dA_0\)，等距保持以零为起点的直线及其初速度，所以 \(A(tv)=b+tQv\)。对平移 \(\tau_\lambda(x)=x+\lambda\)，直接计算
	\[
		A\tau_\lambda A^{-1}=\tau_{Q\lambda}.
	\]
	因此空间形式的等距共轭判据恰变成 \(Q\Lambda=\Lambda'\)。这个条件反过来也使 \(x\mapsto Qx\) 在两个格商上诱导等距。

	取 \(\Lambda=\mathbb Z^2\) 与 \(\Lambda_\varepsilon=\varepsilon\mathbb Z\times\varepsilon^{-1}\mathbb Z\)，其中 \(0<\varepsilon<1\)。两者基本平行四边形面积均为1。前者最短非零格向量长1，后者长 \(\varepsilon\)；正交映射保持格向量长度，所以不可能把两格互相映到，两个面积相同的环面不等距。
\end{solution}

\begin{exercise}\label{g3:ex:even-spherical-spaceforms}
	证明每个连通、完备的 \(2m\) 维正常曲率空间形式（\(m\ge1\)），都等距于相应半径的圆球面或标准实射影空间。
\end{exercise}
\begin{solution}
	由分类定理，它为 \(S^{2m}(a)/\Gamma\)，其中 \(\Gamma\) 自由等距作用。球面等距是环境空间正交变换的限制：在一个球面点固定其像及切微分，便唯一决定一个环境正交映射；二者由大圆测地线的初值唯一性在整个球面相等。因此把 \(\Gamma\) 视为 \(O(2m+1)\) 的子群。

	任意 \(Q\in SO(2m+1)\) 都有特征值1。事实上，非实特征值成共轭对，其乘积为1；剩下奇数个实特征值属于 \(\{1,-1\}\)，总乘积为1，不能全为 \(-1\)。所以非恒等的保定向正交变换若属于 \(\Gamma\)，就会固定一个单位特征向量，与自由性矛盾。于是行列式同态 \(\det:\Gamma\to\{1,-1\}\) 的核平凡，\(|\Gamma|\le2\)。

	若群平凡，商为球面。若群有两个元素，记非恒等元为 \(Q\)，则 \(Q^2=I\)。它的特征值只有 \(\pm1\)，而自由性排除1，故 \(Q=-I\)。所以该商就是标准实射影空间。结论对奇维球面不成立，第八章的 \(S^3\) 循环商便可具有任意给定的有限循环基本群。
\end{solution}

\begin{exercise}\label{g3:ex:parabolic-hyperbolic-cylinder}
	在双曲上半平面上令 \(T(x,y)=(x+1,y)\)。证明 \(\mathbb H^2/\langle T\rangle\) 是完备双曲圆柱，且：
	\begin{enumerate}
		\item 代表一次生成元的闭曲线长度下确界为零，但没有曲线达到该下确界；
		\item 商中 \(y>Y\) 的区域面积为 \(1/Y\)，而整个圆柱面积无限；
		\item 它与正文伸缩 \((x,y)\mapsto(e^\ell x,e^\ell y)\) 所定义的双曲圆柱不等距。
	\end{enumerate}
\end{exercise}
\begin{solution}
	\(T\) 保持度量且其非零次幂无不动点。紧集中 \(x\) 坐标有界，故只有有限个整数平移能使它与自身相交。商作用因而自由、适当不连续，商为完备常曲率 \(-1\) 的曲面。坐标 \(([x],\log y)\) 给出到 \(S^1\times\mathbb R\) 的同胚。

	固定高度 \(y=H\)，曲线 \(x\in[0,1]\mapsto([x],H)\) 代表一次生成元，其长度为 \(1/H\)。令 \(H\to\infty\)，得到下确界零。长度为零的分段光滑曲线只能为常曲线，不能代表非平凡基本群元素，所以这个下确界不达到。

	面积密度为 \(y^{-2}\,\mathrm dx\,\mathrm dy\)。取基本条带 \(0\le x<1\)，有
	\[
		\operatorname{Area}\{y>Y\}
		=\int_0^1\int_Y^\infty y^{-2}\,\mathrm dy\,\mathrm dx=1/Y.
	\]
	将下限令为零时积分发散，故总体积无限。正文伸缩圆柱的任一非平凡闭曲线长度至少为 \(\ell>0\)；这里的非平凡闭曲线长度下确界却为零。等距保持闭曲线长度与同伦的平凡性，故两类圆柱不等距。
\end{solution}

\begin{optionalexercise}\label{g3:ex:toponogov-flat-cylinder-triangle}
 在平坦圆柱 \(M=(\mathbb R/3L\mathbb Z)\times\mathbb R\) 上，\(L>0\)，取
 \[
 p=([0],0),\qquad q=([L],0),\qquad r=([2L],0).
 \]
 三条边分别取圆周上的长度 \(L\) 的短弧。
 \begin{enumerate}
  \item 证明这些边都是最短测地线，三角形的三个角都是 \(\pi\)，而 Euclid 比较三角形的三个角都是 \(\pi/3\)。
  \item 将边 \(qr\) 按 \(\sigma(t)=([L+t],0)\)、\(0\le t\le L\) 参数化，计算 \(\rho(t)=d(p,\sigma(t))\) 和模型距离 \(\bar\rho(t)\)，并验证每个内点都有严格的顶点—对边比较。
  \item 计算 \(h=\rho^2/2\) 在 \(t=L/2\) 的单侧导数，说明本例的严格比较来自哪里。
 \end{enumerate}
\end{optionalexercise}
\begin{solution}
 覆盖平面中的端点提升给出距离公式
 \[
 d(([x],y),([x'],y'))
 =\min_{j\in\mathbb Z}\sqrt{(x'-x+3jL)^2+(y'-y)^2}.
 \]
 三对顶点的距离都为 \(L\)，所以所取短弧都是最短段。在每个顶点，指向另外两个顶点的两条初始方向相反，夹角为 \(\pi\)。模型则是边长 \(L\) 的平面等边三角形，三个角均为 \(\pi/3\)。

 沿所取对边，圆周距离为
 \[
 \rho(t)=\min\{L+t,2L-t\}
 =\begin{cases}L+t,&0\le t\le L/2,\\2L-t,&L/2\le t\le L.\end{cases}
 \]
 在平面比较三角形中，以 \(\bar q=(0,0)\)、\(\bar r=(L,0)\)、\(\bar p=(L/2,\sqrt3L/2)\) 表示顶点，则
 \[
 \bar\rho(t)^2=(t-L/2)^2+3L^2/4=L^2-Lt+t^2.
 \]
 因而
 \[
 \rho(t)^2-\bar\rho(t)^2
 =\begin{cases}3Lt,&0\le t\le L/2,\\3L(L-t),&L/2\le t\le L.\end{cases}
 \]
 两个端点取等号，所有内点严格为正。在中点，实际距离为 \(3L/2\)，模型距离为 \(\sqrt3L/2\)。

 在中点左侧，\(h' = L+t\)；在右侧，\(h'=-(2L-t)\)。因此
 \[
 h'_-(L/2)=3L/2,\qquad h'_+(L/2)=-3L/2.
 \]
 两个光滑段上都有 \(h''=1\)，但在割点处导数向下跳了 \(3L\)。局部曲率虽恒为零，整体最短分支的切换仍使三角比较严格；本例也直接显示一维拼接条件为何允许向下跳。
\end{solution}

\begin{optionalexercise}\label{g3:ex:toponogov-angular-packing}
 设 \((M^n,g)\) 连通、完备、光滑且无边界，\(n\ge2\)、\(K\ge0\)。固定 \(p\in M\)，取不同于 \(p\) 的两两不同点 \(p_1,\ldots,p_N\)，并假设
 \[
 d(p_i,p_j)\ge\max\{d(p,p_i),d(p,p_j)\}\qquad(i\ne j).
 \]
 分别任选从 \(p\) 到 \(p_i\) 的最短段，记其初始单位向量为 \(u_i\)。证明 \(\angle(u_i,u_j)\ge\pi/3\)，并用切空间中的 Euclid 球体积证明 \(N\le3^n\)。这里不要求各个 \(d(p,p_i)\) 相等。
\end{optionalexercise}
\begin{solution}
 记 \(r_i=d(p,p_i)>0\)、\(a_{ij}=d(p_i,p_j)\)。由 Toponogov 角比较，实际夹角不小于具有边长 \(r_i,r_j,a_{ij}\) 的 Euclid 模型角 \(\bar\alpha_{ij}\)。设 \(r_i\le r_j\)，余弦公式与假设给出
 \[
 \cos\bar\alpha_{ij}
 =\frac{r_i^2+r_j^2-a_{ij}^2}{2r_ir_j}
 \le\frac{r_i^2}{2r_ir_j}
 =\frac{r_i}{2r_j}\le\frac12.
 \]
 因而 \(\angle(u_i,u_j)\ge\bar\alpha_{ij}\ge\pi/3\)。这一步对最短段的任意选择均成立，因为角比较本身对每种选择都成立。

 将 \(\mathrm T_pM\) 用内积识别为 \(\mathbb R^n\)。单位向量间的弦长满足
 \[
 |u_i-u_j|^2=2-2\cos\angle(u_i,u_j)\ge1.
 \]
 因此 Euclid 开球 \(B(u_i,1/2)\) 两两不交，且都包含于 \(B(0,3/2)\)。体积相加得到
 \[
 N\omega_n(1/2)^n\le\omega_n(3/2)^n,
 \]
 即 \(N\le3^n\)。所用打包发生在同一个有限维切空间里，所以没有对流形本身的小球体积另加下界。
\end{solution}

\begin{optionalexercise}\label{g3:ex:positive-curvature-equilateral-median}
 设 \((M^n,g)\) 连通、完备、光滑且无边界，\(n\ge2\)、\(K\ge k>0\)。设一个最短测地三角形的三条边长都等于 \(\ell\)，且 \(0<\ell<2b_k/3\)。记 \(z\) 为其中一边的弧长中点，\(p\) 为对顶点，令 \(x=\sqrt{k}\ell\)。证明每个顶角 \(\alpha\) 及中线距离满足
 \[
 \alpha\ge\arccos\frac{\cos x}{1+\cos x},\qquad
 d(p,z)\ge\frac1{\sqrt{k}}\arccos\frac{\cos x}{\cos(x/2)}.
 \]
 核对两个反余弦的自变量范围，并求右侧在 \(\ell\uparrow2b_k/3\) 时的极限。
\end{optionalexercise}
\begin{solution}
 三角周长为 \(3\ell<2b_k\)，所以 Toponogov 的严格周长版本适用。模型角 \(\bar\alpha\) 由球面余弦公式满足
 \[
 \cos x=\cos^2x+\sin^2x\cos\bar\alpha,
 \]
 故 \(\cos\bar\alpha=\cos x/(1+\cos x)\)，角比较给出第一式。在 \(0<x<2\pi/3\) 上有 \(-1/2<\cos x<1\)，从而该商属于 \((-1,1/2)\subset[-1,1]\)。

 为计算模型中线，先缩放到单位球面。设模型顶点在 \(\mathbb R^3\) 中的单位向量为 \(P,Q,R\)，则
 \(\langle P,Q\rangle=\langle P,R\rangle=\langle Q,R\rangle=\cos x\)。由于 \(x<\pi\)，短边 \(QR\) 的中点为
 \[
 Z=\frac{Q+R}{|Q+R|}=\frac{Q+R}{2\cos(x/2)}.
 \]
 所以
 \[
 \langle P,Z\rangle=\frac{\cos x}{\cos(x/2)}.
 \]
 球面内积是距离的余弦，恢复尺度后得到模型中线长度，再用顶点—对边比较即得第二式。

 令 \(y=\cos(x/2)\in(1/2,1)\)，则
 \[
 \frac{\cos x}{\cos(x/2)}=2y-\frac1y.
 \]
 该函数随 \(y\) 严格增加，从 \(-1\) 增至 \(1\)，故本题范围内反余弦合法。当 \(x\uparrow2\pi/3\) 时，两式中的商都趋于 \(-1\)，相应下界分别趋于 \(\pi\) 与 \(\pi/\sqrt{k}=b_k\)。几何上，模型趋向把一条大圆三等分的构形，对边中点趋向对顶点的对径点。
\end{solution}

\begin{exercise}\label{g3:ex:hadamard-barycenter}
	设 \(M\) 是 \(n\ge2\) 维 Hadamard 流形，\(x_1,\ldots,x_N\in M\)，且 \(w_i>0\)、\(\sum_iw_i=1\)。令
	\[
		F(q)=\frac12\sum_{i=1}^Nw_i d(q,x_i)^2.
	\]
	证明 \(F\) 有唯一极小点 \(b\)，它满足 \(\sum_iw_i\log_bx_i=0\)，且
	\[
		\sum_iw_id(q,x_i)^2
		\ge\sum_iw_id(b,x_i)^2+d(b,q)^2.
	\]
	证明每个包含所有 \(x_i\) 的非空闭测地凸集都包含 \(b\)。最后，用一个有限等距变换群的轨道构造这样的 \(F\)，证明该群必有共同不动点；由此推出任一连通完备 \(K\le0\) 流形的基本群无挠。
\end{exercise}
\begin{solution}
	由距离平方 Hessian 估计，\(\operatorname{Hess}F\ge g\)。又因
	\(F(q)\ge(w_1/2)d(q,x_1)^2\)，\(F\) 的有界次水平集是闭有界集，故紧。取一列逼近下确界的点，提取收敛子列，即得极小点 \(b\)。沿不同两点的最短线，二阶导数严格为正，故极小点唯一。梯度公式给出
	\[
		\nabla F(b)=-\sum_iw_i\log_bx_i=0.
	\]
	沿 \(\gamma(t)=\exp_b(t\log_bq)\) 有 \((F\circ\gamma)''\ge d(b,q)^2\) 及 \((F\circ\gamma)'(0)=0\)，积分两次得到
	\(F(q)\ge F(b)+d(b,q)^2/2\)，即所求不等式。反过来，梯度为零的点也由同样论证为唯一极小点。

	若 \(C\) 是题设闭凸集，令 \(u=P_Cb\)。投影的平方距离不等式对每个 \(x_i\in C\) 给出
	\[
		d(b,x_i)^2\ge d(b,u)^2+d(u,x_i)^2.
	\]
	加权求和得到 \(F(b)\ge F(u)+d(b,u)^2/2\)。极小性给出 \(F(b)\le F(u)\)，因而 \(b=u\in C\)。

	设 \(G\) 为有限等距变换群，选定 \(x\in M\)，取
	\[
		F_G(q)=\frac1{2|G|}\sum_{g\in G}d(q,gx)^2.
	\]
	这里允许轨道点重复。对任意 \(h\in G\)，等距性与群元素的置换给出 \(F_G(hq)=F_G(q)\)。它的唯一极小点因此满足 \(hb=b\) 对所有 \(h\in G\) 成立。

	最后取一个连通完备 \(K\le0\) 流形的万有黎曼覆盖。由\cref{g3:prop:riemannian-universal-cover}它完备、单连通且仍满足 \(K\le0\)，所以是 Hadamard 流形；覆盖变换群同构于基本群，且自由作用。如果某个非单位覆盖变换有有限阶，它生成的有限群按刚才的结论有共同不动点，与自由性矛盾。因此基本群没有非平凡有限阶元素。
\end{solution}

\begin{exercise}\label{g3:ex:hyperbolic-projection}
	在双曲上半平面 \(\mathbb H^2=\{(x,y):y>0\}\)、\(g=(\mathrm dx^2+\mathrm dy^2)/y^2\) 中，令 \(C=\{(0,y):y>0\}\)。求点 \((x,y)\) 到 \(C\) 的最近点和距离。可从曲线
	\[
		\gamma(t)=(r\tanh t,r/\cosh t)
	\]
	出发，验证它是单位速测地线，并使用最近点的角刻画。比较所得投影与 Euclid 平面到竖直线的正交投影。
\end{exercise}
\begin{solution}
	由度量的坐标公式，非零 Christoffel 符号为
	\[
		\Gamma^x_{xy}=\Gamma^x_{yx}=-1/y,\qquad
		\Gamma^y_{xx}=1/y,\qquad \Gamma^y_{yy}=-1/y.
	\]
	故测地线方程为
	\(x''-2x'y'/y=0\)、\(y''+((x')^2-(y')^2)/y=0\)。
	对给定曲线，
	\[
		x'=r\,\operatorname{sech}^2t,\qquad
		y'=-r\,\operatorname{sech}t\tanh t,
	\]
	代入两条方程均得零，并由
	\((x')^2+(y')^2=r^2\operatorname{sech}^2t=y^2\)
	得到单位速度。

	取 \(r=\sqrt{x^2+y^2}\) 及 \(t_0=\operatorname{arsinh}(x/y)\)，就有 \(\gamma(t_0)=(x,y)\)、\(\gamma(0)=(0,r)\)。竖直线以 \((0,e^s)\) 参数化也是测地线；在 Hadamard 双曲模型中它是闭测地凸集。曲线 \(\gamma\) 在 \(t=0\) 的切向为水平，因而与 \(C\) 垂直；该度量与 Euclid 度量共形，所以直角的判断相同。对 \(C\) 上所有 \(q\)，向量 \(\log_{(0,r)}q\) 都沿竖直方向，满足投影角条件。因此
	\[
		P_C(x,y)=(0,\sqrt{x^2+y^2}),\qquad
		d((x,y),C)=|t_0|
		=\operatorname{arsinh}\frac{|x|}{y}.
	\]
	这里单位速测地线的两点间距等于参数差，使用了双曲模型的全局唯一最短线性质。Euclid 投影则为 \((0,y)\)，只有 \(x=0\) 时二者相同。相同的角度概念配上不同的测地线，会给出不同的垂足。
\end{solution}

\begin{exercise}\label{g3:ex:cat-zero-quadrilateral}
	设 \(M\) 是 \(n\ge2\) 维 Hadamard 流形，\(a,b,c,d\in M\)。证明
	\[
		d(a,c)^2+d(b,d)^2
		\le d(a,b)^2+d(b,c)^2+d(c,d)^2+d(d,a)^2.
	\]
	在 Euclid 空间中，求右端减去左端的精确表达式，并刻画等号。
\end{exercise}
\begin{solution}
	令 \(m\) 为 \([a,c]\) 的中点。式\eqref{g3:eq:hadamard-cn}分别给出
	\[
	\begin{aligned}
		d(b,m)^2&\le\tfrac12d(b,a)^2+\tfrac12d(b,c)^2-\tfrac14d(a,c)^2,\\
		d(d,m)^2&\le\tfrac12d(d,a)^2+\tfrac12d(d,c)^2-\tfrac14d(a,c)^2.
	\end{aligned}
	\]
	另一方面，三角不等式和 \((u+v)^2\le2u^2+2v^2\) 给出
	\[
		d(b,d)^2\le2d(b,m)^2+2d(d,m)^2.
	\]
	代入前两式并整理，就得到所求四边形不等式。

	若各点为 Euclid 向量，直接展开四个边长平方并减去两个对角线平方，得到
	\[
		|a-b|^2+|b-c|^2+|c-d|^2+|d-a|^2-|a-c|^2-|b-d|^2
		=|a+c-b-d|^2.
	\]
	故等号当且仅当 \(a+c=b+d\)，即两条对角线具有相同中点；这正是允许退化情形的平行四边形条件。
\end{solution}
